\chapter{Heaps}
\chapterepigraph{A heap does not know the full order of its elements, and
it does not need to. It only ever promises one thing: the element at the
root is the best (smallest, or largest) of everything currently inside --
and that single promise, maintained cheaply after every insertion and
removal, is enough to power an entire family of ``best/kth/top'' problems.}

\section{What Is a Heap, Really?}
\label{sec:heap-intro}

A \textbf{heap} (also called a \textbf{priority queue} when discussed in
terms of its interface rather than its array-based implementation) is a
tree-shaped data structure, almost always stored compactly in a plain
array, that maintains one structural invariant: in a \textbf{min-heap},
every node's value is $\le$ the values of its children; in a
\textbf{max-heap}, every node's value is $\ge$ the values of its children.
This invariant is weaker than full sortedness -- a heap does not know the
relative order of two sibling subtrees -- but it is exactly strong enough
to guarantee that the root is always the minimum (or maximum) of the whole
structure, and that this root can be removed and the invariant restored in
$O(\log n)$ time.

\begin{ideabox}[The Three Operations That Power Everything in This Chapter]
\begin{itemize}
  \item $\textit{push}(x)$: insert $x$ at the next free array slot, then
        \emph{sift up} -- repeatedly swap it with its parent while it
        violates the heap invariant relative to that parent. Cost:
        $O(\log n)$.
  \item $\textit{pop}()$: read the root (the answer), move the last
        element into the root's slot, shrink the array by one, then
        \emph{sift down} -- repeatedly swap the new root with whichever
        child violates the invariant most, until the invariant holds
        again. Cost: $O(\log n)$.
  \item $\textit{top}()$: read the root without removing it. Cost: $O(1)$.
\end{itemize}
Every problem in this chapter is, underneath, a creative use of just these
three operations.
\end{ideabox}

\subsection*{How to Recognise a Heap Problem from Its Wording}

\begin{patternbox}[Recurring Heap Keywords]
\begin{itemize}
  \item \textbf{``Kth largest/smallest''} -- maintain a heap of size $k$
        of the opposite polarity (a min-heap of size $k$ for ``$k$th
        largest,'' so the root is the worst of your current top-$k$, ready
        to be evicted the instant something better arrives).
  \item \textbf{``Top $k$ frequent/most/best''} -- the same size-$k$ heap
        trick, ranked by a derived key (frequency, sum, distance) rather
        than the raw value.
  \item \textbf{``Merge $k$ sorted ...''} -- a heap holding one
        ``current head'' candidate from each of the $k$ sources at a time,
        repeatedly popping the smallest and pushing that source's next
        element.
  \item \textbf{``Running/streaming median, or any other running
        statistic that needs both a top half and a bottom half''} -- two
        heaps, one min-heap and one max-heap, used together as a
        self-balancing pair.
  \item \textbf{``Schedule/group/connect with a cost or cooldown that
        always prefers the largest (or smallest) remaining item first''}
        -- a single heap used as a perpetually-resorted greedy priority
        queue, replacing what would otherwise require re-sorting after
        every step.
\end{itemize}
Whenever brute force would require repeatedly finding ``the current
best/worst remaining element'' after the underlying collection keeps
changing, a heap is almost always the efficient replacement for ``sort
again from scratch every time.''
\end{patternbox}

Each section in this chapter follows the same structure as the rest of
this book: the problem statement and its keywords; the reasoning that
identifies which heap (or pair of heaps) to use and why; the algorithm in
words and in formal pseudocode; a hand-traced dry run; complete compilable
C++ code (using \texttt{std::priority\_queue} throughout, since that is
what production C++ code -- and nearly every interview answer -- actually
uses); and a time/space complexity analysis.

\newpage

\section{Implement a Min Heap}
\label{sec:heap-implement}

\problemstatement{Implement a min-heap supporting three operations:
$\textit{push}(x)$ (insert a new value), $\textit{pop}()$ (remove and
return the smallest value), and $\textit{top}()$ (peek at the smallest
value without removing it), each in $O(\log n)$ time for push/pop and
$O(1)$ for top.}

\begin{patternbox}
This section has no ``keyword'' in the usual sense -- it is the mechanical
foundation every other section in this chapter builds on. The skill being
tested is not pattern recognition but the array-indexing arithmetic that
lets a heap live inside a flat array without ever needing explicit
left/right/parent pointers.
\end{patternbox}

\subsection*{Understanding the array-as-tree representation}

Store the heap's elements in a plain array $\textit{heap}[0..n-1]$,
interpreted as a complete binary tree read level by level, left to right.
For a node at index $i$ (0-indexed):
\[
  \textit{parent}(i) = \left\lfloor \frac{i-1}{2} \right\rfloor, \qquad
  \textit{left}(i) = 2i+1, \qquad \textit{right}(i) = 2i+2.
\]
This works because a complete binary tree (every level full except
possibly the last, which fills left to right with no gaps) has exactly
this relationship between a node's position in level order and its
children's positions -- no pointers are needed, only arithmetic.

\begin{ideabox}[Sift Up (used by push)]
Append the new value at index $n$ (the next free slot). While this node's
value is smaller than its parent's value, swap them and continue checking
from the parent's old position. This restores the min-heap invariant
locally along the single path from the new leaf up to the root, which is
always sufficient: every other path in the tree was already valid before
the insertion and is untouched by these swaps.
\end{ideabox}

\begin{ideabox}[Sift Down (used by pop)]
Save the root's value (this is the return value). Move the \emph{last}
element into the root's slot and shrink the array by one. While this
relocated value is larger than the smaller of its two children, swap it
with that smaller child and continue checking from the new position. This
restores the invariant locally along a single path from the root downward.
\end{ideabox}

\subsection*{Why sifting along a single path is always sufficient}

Both operations rely on the same structural fact: the min-heap invariant
is a \emph{local} condition (each node versus its own parent/children),
and inserting or removing one element only disturbs the invariant along
\emph{one} root-to-leaf path -- the path the new or relocated element
travels. Every other node's relationship to its own parent and children is
completely unaffected, so checking and fixing violations one step at a
time along that single path, stopping the instant no violation remains, is
guaranteed to restore a fully valid heap. This is also exactly why each
operation costs $O(\log n)$: a path from root to leaf in a complete binary
tree with $n$ nodes has length $O(\log n)$, since the tree's height is
$\lfloor \log_2 n \rfloor$.

\subsection*{Algorithm}

\textbf{Push($x$):}
\begin{enumerate}
  \item Append $x$ at index $n$; let $i \gets n$; increment $n$.
  \item While $i > 0$ and $\textit{heap}[\textit{parent}(i)] >
        \textit{heap}[i]$: swap them; $i \gets \textit{parent}(i)$.
\end{enumerate}

\textbf{Pop():}
\begin{enumerate}
  \item Save $\textit{result} \gets \textit{heap}[0]$.
  \item $\textit{heap}[0] \gets \textit{heap}[n-1]$; decrement $n$.
  \item $i \gets 0$. While $i$ has a child smaller than itself: let
        $\textit{smallerChild}$ be the smaller of its two children (only
        those that exist); swap $\textit{heap}[i]$ with
        $\textit{heap}[\textit{smallerChild}]$; $i \gets
        \textit{smallerChild}$.
  \item Return \textit{result}.
\end{enumerate}

\begin{algorithm}[H]
\caption{Min Heap: Push and Pop}
\begin{algorithmic}[1]
\Procedure{Push}{$x$}
  \State append $x$; $i \gets n - 1$
  \While{$i > 0$ \textbf{and} $\textit{heap}[(i-1)/2] > \textit{heap}[i]$}
    \State swap $\textit{heap}[(i-1)/2]$ and $\textit{heap}[i]$
    \State $i \gets (i-1)/2$
  \EndWhile
\EndProcedure
\Procedure{Pop}{}
  \State $\textit{result} \gets \textit{heap}[0]$
  \State $\textit{heap}[0] \gets \textit{heap}[n-1]$; remove last slot
  \State $i \gets 0$
  \While{$i$ has a smaller child}
    \State $j \gets$ index of the smaller existing child of $i$
    \If{$\textit{heap}[i] \le \textit{heap}[j]$} \textbf{break} \EndIf
    \State swap $\textit{heap}[i]$ and $\textit{heap}[j]$; $i \gets j$
  \EndWhile
  \State \Return \textit{result}
\EndProcedure
\end{algorithmic}
\end{algorithm}

\subsection*{Dry run}

\dryrun{Push $5, 3, 8, 1$ in order, then pop once.}

\begin{itemize}
  \item Push $5$: heap $=[5]$.
  \item Push $3$: append, heap $=[5,3]$. Compare $3$ to parent $5$:
        $3<5$, swap. Heap $=[3,5]$.
  \item Push $8$: append, heap $=[3,5,8]$. Parent of index $2$ is index
        $0$, value $3$. $8 \ge 3$, no swap. Heap $=[3,5,8]$.
  \item Push $1$: append, heap $=[3,5,8,1]$. Parent of index $3$ is index
        $1$, value $5$. $1<5$, swap: heap $=[3,1,8,5]$. New index $1$;
        parent is index $0$, value $3$. $1<3$, swap: heap $=[1,3,8,5]$.
  \item Pop: result $=1$. Move last element ($5$) to root: heap
        $=[5,3,8]$. At index $0$ ($5$): children at indices $1,2$ are
        $3,8$; smaller is $3$ at index $1$. $5>3$, swap: heap
        $=[3,5,8]$. At index $1$ ($5$): only possible child would be at
        index $3$, out of range. Stop.
\end{itemize}

The pop returns $1$ (correctly the minimum), and the heap array
$[3,5,8]$ still satisfies the min-heap invariant.

\subsection*{C++ implementation}

\begin{lstlisting}
#include <bits/stdc++.h>
using namespace std;

class MinHeap {
    vector<int> heap;

    void siftUp(int i) {
        while (i > 0 && heap[(i - 1) / 2] > heap[i]) {
            swap(heap[(i - 1) / 2], heap[i]);
            i = (i - 1) / 2;
        }
    }

    void siftDown(int i) {
        int n = heap.size();
        while (true) {
            int left = 2 * i + 1, right = 2 * i + 2, smallest = i;
            if (left < n && heap[left] < heap[smallest]) smallest = left;
            if (right < n && heap[right] < heap[smallest]) smallest = right;
            if (smallest == i) break;
            swap(heap[i], heap[smallest]);
            i = smallest;
        }
    }

public:
    void push(int x) {
        heap.push_back(x);
        siftUp(heap.size() - 1);
    }

    int pop() {
        int result = heap[0];
        heap[0] = heap.back();
        heap.pop_back();
        if (!heap.empty()) siftDown(0);
        return result;
    }

    int top() { return heap[0]; }
    bool empty() { return heap.empty(); }
};

int main() {
    MinHeap h;
    h.push(5); h.push(3); h.push(8); h.push(1);
    cout << "Popped: " << h.pop() << endl;
    cout << "New top: " << h.top() << endl;
    return 0;
}
\end{lstlisting}

\begin{complexitybox}
\textbf{Time Complexity.} Both \textit{push} and \textit{pop} touch only
one root-to-leaf path, of length $O(\log n)$, doing $O(1)$ work per level,
giving
\[
  O(\log n) \text{ per push/pop}, \qquad O(1) \text{ per top}.
\]
\textbf{Space Complexity.} $O(n)$ to store the $n$ elements; no extra
auxiliary structures are needed beyond the array itself.
\end{complexitybox}

\begin{takeawaybox}
Every problem in the rest of this chapter uses \texttt{std::priority\_queue}
directly rather than reimplementing this class, but understanding sift-up
and sift-down by hand here is what makes the complexity claims in every
later section ($O(\log n)$ per heap operation) something you can justify
rather than merely recite.
\end{takeawaybox}

\section{Check if an Array Represents a Heap}
\label{sec:heap-check-array}

\problemstatement{Given an array $\textit{arr}[0..n-1]$, determine whether
it represents a valid \textbf{max-heap} when interpreted as a complete
binary tree in level order (every parent at index $i$ should be $\ge$ its
children at $2i+1$ and $2i+2$).}

\begin{patternbox}
This is a direct test of the array-as-tree arithmetic from
Section~\ref{sec:heap-implement}: the keyword is simply \emph{``does this
array, read as a complete binary tree, satisfy the heap invariant
everywhere.''} There is no clever algorithm to discover -- the entire
exercise is checking the invariant at every internal node using the
$\textit{left}(i) = 2i+1$, $\textit{right}(i) = 2i+2$ formulas directly.
\end{patternbox}

\subsection*{Understanding the problem carefully}

A leaf node (one with no children, i.e.\ $2i+1 \ge n$) trivially satisfies
the heap property, since it has nothing to compare against. So we only need
to check internal nodes -- indices $i$ with $2i+1 < n$ -- and verify that
$\textit{arr}[i] \ge \textit{arr}[2i+1]$ (when the left child exists) and
$\textit{arr}[i] \ge \textit{arr}[2i+2]$ (when the right child exists). If
every internal node passes both checks, the whole array is a valid
max-heap, because the heap property is purely local: validity at every
parent-child pair, checked independently, is both necessary and sufficient
for the global invariant to hold (there is no broader ordering constraint
a heap is required to satisfy beyond these local pairwise comparisons).

\begin{ideabox}[Why Checking Only Internal Nodes Suffices]
The heap property is defined entirely in terms of parent-versus-children
comparisons. A node with no children imposes no constraint of its own, so
the full set of constraints to verify is exactly the set of (parent, child)
pairs that exist in the array -- checking every internal node against its
existing children covers all of them exactly once.
\end{ideabox}

\subsection*{Algorithm}

\begin{enumerate}
  \item For $i \gets 0$ to $n-1$, while $2i+1 < n$ (i.e.\ $i$ is internal):
  \begin{enumerate}
    \item If $2i+1 < n$ and $\textit{arr}[i] < \textit{arr}[2i+1]$: return
          \texttt{false}.
    \item If $2i+2 < n$ and $\textit{arr}[i] < \textit{arr}[2i+2]$: return
          \texttt{false}.
  \end{enumerate}
  \item If every internal node passes, return \texttt{true}.
\end{enumerate}

\begin{algorithm}[H]
\caption{Check Array Represents a Max-Heap}
\begin{algorithmic}[1]
\Procedure{IsHeap}{\textit{arr}}
  \For{$i \gets 0$ to $n-1$}
    \State $\textit{left} \gets 2i+1$, $\textit{right} \gets 2i+2$
    \If{$\textit{left} < n$ \textbf{and} $\textit{arr}[i] < \textit{arr}[\textit{left}]$}
      \State \Return \textbf{false}
    \EndIf
    \If{$\textit{right} < n$ \textbf{and} $\textit{arr}[i] < \textit{arr}[\textit{right}]$}
      \State \Return \textbf{false}
    \EndIf
  \EndFor
  \State \Return \textbf{true}
\EndProcedure
\end{algorithmic}
\end{algorithm}

\subsection*{Dry run}

\dryrun{Take $\textit{arr} = [90, 15, 10, 7, 12, 2, 7, 3]$.}

\begin{itemize}
  \item $i=0$ (value $90$): children at $1,2$ are $15,10$. $90 \ge 15$ and
        $90 \ge 10$. OK.
  \item $i=1$ (value $15$): children at $3,4$ are $7,12$. $15 \ge 7$ and
        $15 \ge 12$. OK.
  \item $i=2$ (value $10$): children at $5,6$ are $2,7$. $10 \ge 2$ and
        $10 \ge 7$. OK.
  \item $i=3$ (value $7$): child at $7$ is $3$ (no right child, index $8$
        out of range). $7 \ge 3$. OK.
  \item $i=4,5,6,7$: no children ($2i+1 \ge 8$). Trivially OK.
\end{itemize}

Every internal node passes, so the array is a valid max-heap.

\subsection*{C++ implementation}

\begin{lstlisting}
#include <bits/stdc++.h>
using namespace std;

bool isHeap(vector<int>& arr) {
    int n = arr.size();
    for (int i = 0; i < n; i++) {
        int left = 2 * i + 1, right = 2 * i + 2;
        if (left < n && arr[i] < arr[left]) return false;
        if (right < n && arr[i] < arr[right]) return false;
    }
    return true;
}

int main() {
    vector<int> arr = {90, 15, 10, 7, 12, 2, 7, 3};
    cout << (isHeap(arr) ? "true" : "false") << endl;
    return 0;
}
\end{lstlisting}

\begin{complexitybox}
\textbf{Time Complexity.} Each index is visited once with constant work,
giving
\[
  O(n).
\]
\textbf{Space Complexity.} $O(1)$ auxiliary space.
\end{complexitybox}

\begin{takeawaybox}
The heap property is purely local (parent versus its own children), never
global -- a heap is not required to be sorted in any stronger sense. This
is precisely why heap operations are $O(\log n)$ rather than $O(n \log n)$:
restoring a single local violation after an insertion or removal never
requires re-examining the whole structure.
\end{takeawaybox}

\section{Convert a Min Heap to a Max Heap}
\label{sec:heap-min-to-max}

\problemstatement{Given an array that currently satisfies the min-heap
property, rearrange it in place so that it satisfies the max-heap property
instead.}

\begin{patternbox}
The keyword here is \emph{``convert,'' not ``rebuild from a sorted
list''} -- the existing min-heap order is, in fact, almost entirely
irrelevant to solving this efficiently; the fastest approach simply
\textbf{re-heapifies the array from scratch as a max-heap}, ignoring its
current (min-heap) structure. This problem is the natural place to
introduce the classic $O(n)$ \textbf{build-heap} algorithm, which is
faster than the $O(n \log n)$ you would get by popping every element from
the min-heap and pushing it into a fresh max-heap one at a time.
\end{patternbox}

\subsection*{Understanding the problem carefully}

The simplest correct approach: treat the existing array's \emph{values}
as an unordered list of $n$ numbers (discard the fact that they currently
satisfy a min-heap arrangement -- it buys us nothing for the new task),
and re-heapify in place into a max-heap. The key efficiency insight is
that \emph{any} array can be turned into a valid heap in $O(n)$ time, not
$O(n \log n)$, by sifting down starting from the \textbf{last internal
node} and working backward to the root.

\begin{ideabox}[Why Starting from the Last Internal Node Works]
A leaf node trivially satisfies the heap property on its own. If we
process internal nodes in \emph{reverse} level order (bottom-most,
right-most internal node first, ending at the root), then by the time we
sift-down a given node $i$, both of its children's subtrees have
\emph{already} been fixed into valid max-heaps (since they were processed
earlier, being lower in the tree). Sifting node $i$ down then only needs
to resolve a violation between $i$ and the roots of two already-valid
subtrees -- exactly the situation sift-down is designed to fix -- so a
single bottom-up pass, sifting each internal node down once, is enough to
turn the entire array into a valid heap.
\end{ideabox}

\subsection*{Why this build-heap process is $O(n)$, not $O(n \log n)$}

Each sift-down costs time proportional to the \emph{height of the subtree
rooted at that node}, not the height of the whole tree. Most nodes are
near the bottom of the tree, where subtree heights are tiny; only the
handful of nodes near the root have a large subtree height. Formally, the
total work is
\[
  \sum_{h=0}^{\log n} \frac{n}{2^{h+1}} \cdot h = O(n),
\]
since the number of nodes at height $h$ from the bottom is roughly
$n/2^{h+1}$, and this geometric-times-linear sum converges to $O(n)$
rather than the naive (and incorrect) $O(n \log n)$ upper bound you would
get by assuming every node costs the full tree height.

\subsection*{Algorithm}

\begin{enumerate}
  \item The last internal node (one with at least one child) is at index
        $\lfloor n/2 \rfloor - 1$.
  \item For $i$ from $\lfloor n/2 \rfloor - 1$ down to $0$: sift-down
        node $i$, using max-heap comparisons (swap with the \emph{larger}
        child whenever a child exceeds the current node).
  \item The array now satisfies the max-heap property.
\end{enumerate}

\begin{algorithm}[H]
\caption{Build a Max-Heap In Place (Heapify)}
\begin{algorithmic}[1]
\Procedure{BuildMaxHeap}{\textit{arr}}
  \For{$i \gets \lfloor n/2 \rfloor - 1$ \textbf{down to} $0$}
    \State \Call{SiftDownMax}{\textit{arr}, $i$, $n$}
  \EndFor
\EndProcedure
\Procedure{SiftDownMax}{\textit{arr}, $i$, $n$}
  \While{\textbf{true}}
    \State $\textit{largest} \gets i$, $\textit{left} \gets 2i+1$, $\textit{right} \gets 2i+2$
    \If{$\textit{left} < n$ \textbf{and} $\textit{arr}[\textit{left}] > \textit{arr}[\textit{largest}]$}
      \State $\textit{largest} \gets \textit{left}$
    \EndIf
    \If{$\textit{right} < n$ \textbf{and} $\textit{arr}[\textit{right}] > \textit{arr}[\textit{largest}]$}
      \State $\textit{largest} \gets \textit{right}$
    \EndIf
    \If{$\textit{largest} = i$} \textbf{break} \EndIf
    \State swap $\textit{arr}[i]$ and $\textit{arr}[\textit{largest}]$; $i \gets \textit{largest}$
  \EndWhile
\EndProcedure
\end{algorithmic}
\end{algorithm}

\subsection*{Dry run}

\dryrun{Take the min-heap $\textit{arr} = [1,3,5,4,6,7,8]$ ($n=7$).}

Last internal node: $\lfloor 7/2 \rfloor - 1 = 2$.

\begin{itemize}
  \item $i=2$ (value $5$): children at $5,6$ are $7,8$. Largest child is
        $8$ at index $6$. $8>5$: swap. $\textit{arr}=[1,3,8,4,6,7,5]$.
        Index $6$ is a leaf; stop.
  \item $i=1$ (value $3$): children at $3,4$ are $4,6$. Largest is $6$ at
        index $4$. $6>3$: swap. $\textit{arr}=[1,6,8,4,3,7,5]$. Index $4$
        is a leaf; stop.
  \item $i=0$ (value $1$): children at $1,2$ are $6,8$. Largest is $8$ at
        index $2$. $8>1$: swap. $\textit{arr}=[8,6,1,4,3,7,5]$. Now at
        index $2$ (value $1$): children at $5,6$ are $7,5$. Largest is
        $7$ at index $5$. $7>1$: swap.
        $\textit{arr}=[8,6,7,4,3,1,5]$. Index $5$ is a leaf; stop.
\end{itemize}

The final array $[8,6,7,4,3,1,5]$ satisfies the max-heap property (the
root $8$ is the maximum, and every parent exceeds its children).

\subsection*{C++ implementation}

\begin{lstlisting}
#include <bits/stdc++.h>
using namespace std;

void siftDownMax(vector<int>& arr, int i, int n) {
    while (true) {
        int largest = i, left = 2 * i + 1, right = 2 * i + 2;
        if (left < n && arr[left] > arr[largest]) largest = left;
        if (right < n && arr[right] > arr[largest]) largest = right;
        if (largest == i) break;
        swap(arr[i], arr[largest]);
        i = largest;
    }
}

void minHeapToMaxHeap(vector<int>& arr) {
    int n = arr.size();
    for (int i = n / 2 - 1; i >= 0; i--) {
        siftDownMax(arr, i, n);
    }
}

int main() {
    vector<int> arr = {1, 3, 5, 4, 6, 7, 8}; // valid min-heap
    minHeapToMaxHeap(arr);

    cout << "Max-heap: ";
    for (int x : arr) cout << x << " ";
    cout << endl;
    return 0;
}
\end{lstlisting}

\begin{complexitybox}
\textbf{Time Complexity.} The build-heap process visits $O(n/2)$ internal
nodes with total sift-down work summing to
\[
  O(n)
\]
across the whole array (not $O(n \log n)$), by the height-weighted
geometric sum argument above.
\textbf{Space Complexity.} $O(1)$ auxiliary space (the conversion is done
in place).
\end{complexitybox}

\begin{takeawaybox}
Re-heapifying an entire array from scratch is $O(n)$, not $O(n \log n)$ --
a fact that is easy to misremember because a \emph{single} sift-down is
$O(\log n)$, tempting you to multiply by $n$ insertions. The $O(n)$ bound
only holds when you build the whole heap at once via bottom-up sift-down,
not when you insert $n$ elements one at a time via repeated sift-up.
\textsc{std::make\_heap} in C++ uses exactly this $O(n)$ bottom-up
technique internally.
\end{takeawaybox}

\section{Kth Largest Element in an Array}
\label{sec:heap-kth-largest}

\problemstatement{Given an array $\textit{nums}$ of $n$ integers and an
integer $k$, find the $k$-th largest element in the array (the element
that would be at index $n-k$ if the array were sorted in increasing
order). Note this is the $k$-th largest \emph{value}, counting duplicates
by position, not the $k$-th distinct value.}

\begin{patternbox}
\emph{``Kth largest''} is the single most direct trigger for the
size-bounded heap technique introduced in
Section~\ref{sec:heap-intro}'s keyword list. The key (and slightly
counter-intuitive) design choice: to find the $k$-th \textbf{largest}
element efficiently, maintain a \textbf{min}-heap, not a max-heap -- the
min-heap's root is the \emph{worst} element among your current best-$k$
candidates, which is exactly the one you want instant access to evict.
\end{patternbox}

\subsection*{Understanding the problem carefully}

Sorting the whole array and reading off index $n-k$ costs
$O(n \log n)$ and uses the full order of all $n$ elements, even though we
only care about one specific rank. A heap lets us track only the $k$
largest elements seen so far, discarding everything else immediately.

\begin{ideabox}[Why a Min-Heap, not a Max-Heap]
Maintain a min-heap that never holds more than $k$ elements -- think of it
as ``the best $k$ candidates for being among the $k$ largest, with the
\emph{worst} of those $k$ candidates sitting at the root, ready to be
evicted.'' Process the array one element at a time: push the new element,
then if the heap's size exceeds $k$, pop the root (the smallest among the
current top-$k$ candidates, which is therefore the easiest one to safely
discard). After processing every element, the heap holds exactly the $k$
largest elements of the whole array, and its root -- the smallest of
those $k$ -- is precisely the $k$-th largest element overall.
\end{ideabox}

\subsection*{Why evicting the heap's root is always a safe discard}

At any point during the scan, the min-heap holds some $k$ candidates for
``currently in the top $k$.'' When a new element arrives and the heap
already has $k$ elements, we must decide: is the new element large enough
to belong in the top $k$, displacing the current weakest member? Pushing
the new element first and then popping the minimum handles both cases
uniformly without an explicit if-statement: if the new element is smaller
than everything already in the heap, it becomes the new minimum and is
immediately popped right back out, leaving the heap unchanged in effect;
if the new element is larger than the current minimum, that old minimum
(now provably not among the true top $k$, since at least $k$ other
elements, including the new one, are now known to be at least as large)
is correctly evicted instead. Either way, the heap always ends each step
holding the $k$ largest elements seen \emph{so far}, by induction.

\subsection*{Algorithm}

\begin{enumerate}
  \item Initialize an empty min-heap.
  \item For each element $x$ in $\textit{nums}$:
  \begin{enumerate}
    \item Push $x$ onto the heap.
    \item If the heap's size exceeds $k$: pop the minimum.
  \end{enumerate}
  \item Return the heap's current minimum (its root).
\end{enumerate}

\begin{algorithm}[H]
\caption{Kth Largest Element}
\begin{algorithmic}[1]
\Procedure{KthLargest}{\textit{nums}, $k$}
  \State $\textit{heap} \gets$ empty min-heap
  \For{each $x$ in \textit{nums}}
    \State push $x$ onto \textit{heap}
    \If{$|\textit{heap}| > k$}
      \State pop the minimum from \textit{heap}
    \EndIf
  \EndFor
  \State \Return top of \textit{heap}
\EndProcedure
\end{algorithmic}
\end{algorithm}

\subsection*{Dry run}

\dryrun{Take $\textit{nums} = [3,2,1,5,6,4]$, $k = 2$.}

\begin{itemize}
  \item Push $3$: heap $=\{3\}$. Size $1 \le 2$.
  \item Push $2$: heap $=\{2,3\}$. Size $2 \le 2$.
  \item Push $1$: heap $=\{1,2,3\}$. Size $3>2$: pop min ($1$). Heap
        $=\{2,3\}$.
  \item Push $5$: heap $=\{2,3,5\}$. Size $3>2$: pop min ($2$). Heap
        $=\{3,5\}$.
  \item Push $6$: heap $=\{3,5,6\}$. Size $3>2$: pop min ($3$). Heap
        $=\{5,6\}$.
  \item Push $4$: heap $=\{4,5,6\}$. Size $3>2$: pop min ($4$). Heap
        $=\{5,6\}$.
\end{itemize}

The final heap holds $\{5,6\}$, the two largest elements; its root
(minimum of these two) is $5$, the $2$nd largest element overall. Sorting
$[1,2,3,4,5,6]$ confirms index $n-k=4$ holds $5$.

\subsection*{C++ implementation}

\begin{lstlisting}
#include <bits/stdc++.h>
using namespace std;

int kthLargest(vector<int>& nums, int k) {
    priority_queue<int, vector<int>, greater<int>> minHeap; // min-heap

    for (int x : nums) {
        minHeap.push(x);
        if ((int)minHeap.size() > k) {
            minHeap.pop();
        }
    }
    return minHeap.top();
}

int main() {
    vector<int> nums = {3, 2, 1, 5, 6, 4};
    cout << "2nd largest: " << kthLargest(nums, 2) << endl;
    return 0;
}
\end{lstlisting}

\begin{complexitybox}
\textbf{Time Complexity.} Each of the $n$ elements triggers one push, and
at most one pop, each costing $O(\log k)$ since the heap never exceeds
size $k$, giving an overall time complexity of
\[
  O(n \log k),
\]
a clear improvement over $O(n \log n)$ when $k \ll n$.
\textbf{Space Complexity.} $O(k)$, since the heap never holds more than
$k$ elements.
\end{complexitybox}

\begin{takeawaybox}
The ``opposite polarity'' rule is worth memorizing precisely because it
feels backward at first: for \emph{largest}-$k$ problems, use a
\emph{min}-heap of size $k$ (so the easiest-to-discard candidate is always
at the root); for \emph{smallest}-$k$ problems
(Section~\ref{sec:heap-kth-smallest}), use a \emph{max}-heap of size $k$.
The heap's root is always the \emph{boundary} element of your current
best-$k$ set -- the one you're most willing to give up.
\end{takeawaybox}

\section{Kth Smallest Element in an Array}
\label{sec:heap-kth-smallest}

\problemstatement{Given an array $\textit{nums}$ of $n$ integers and an
integer $k$, find the $k$-th smallest element in the array.}

\begin{patternbox}
The mirror image of Section~\ref{sec:heap-kth-largest}: \emph{``Kth
smallest''} calls for the \emph{opposite}-polarity heap, a
\textbf{max}-heap of size $k$. The root of a size-$k$ max-heap is the
\emph{largest} of the current best-$k$-smallest candidates -- exactly the
one to discard the instant a smaller element arrives.
\end{patternbox}

\subsection*{Understanding the problem carefully}

By the same reasoning as Section~\ref{sec:heap-kth-largest}, applied with
every comparison flipped: maintain a max-heap that never holds more than
$k$ elements. Push each new element, then if the heap exceeds size $k$,
pop the maximum (the current weakest member of the ``smallest $k$
so far'' set). After processing the whole array, the heap holds the $k$
smallest elements, and its root -- the largest of those $k$ -- is the
$k$-th smallest element overall.

\begin{ideabox}[Reuse, with Every Comparison Flipped]
Apply Algorithm~\ref{sec:heap-kth-largest}'s \textsc{KthLargest} procedure
verbatim, replacing the min-heap with a max-heap and ``pop the minimum''
with ``pop the maximum.''
\end{ideabox}

\subsection*{Algorithm}

\begin{enumerate}
  \item Initialize an empty max-heap.
  \item For each element $x$ in $\textit{nums}$:
  \begin{enumerate}
    \item Push $x$ onto the heap.
    \item If the heap's size exceeds $k$: pop the maximum.
  \end{enumerate}
  \item Return the heap's current maximum (its root).
\end{enumerate}

\begin{algorithm}[H]
\caption{Kth Smallest Element}
\begin{algorithmic}[1]
\Procedure{KthSmallest}{\textit{nums}, $k$}
  \State $\textit{heap} \gets$ empty max-heap
  \For{each $x$ in \textit{nums}}
    \State push $x$ onto \textit{heap}
    \If{$|\textit{heap}| > k$}
      \State pop the maximum from \textit{heap}
    \EndIf
  \EndFor
  \State \Return top of \textit{heap}
\EndProcedure
\end{algorithmic}
\end{algorithm}

\subsection*{Dry run}

\dryrun{Take $\textit{nums} = [7,10,4,3,20,15]$, $k = 3$.}

\begin{itemize}
  \item Push $7$: heap $=\{7\}$.
  \item Push $10$: heap $=\{7,10\}$.
  \item Push $4$: heap $=\{4,7,10\}$. Size $3 \le 3$.
  \item Push $3$: heap $=\{3,4,7,10\}$. Size $4>3$: pop max ($10$). Heap
        $=\{3,4,7\}$.
  \item Push $20$: heap $=\{3,4,7,20\}$. Size $4>3$: pop max ($20$). Heap
        $=\{3,4,7\}$.
  \item Push $15$: heap $=\{3,4,7,15\}$. Size $4>3$: pop max ($15$). Heap
        $=\{3,4,7\}$.
\end{itemize}

The final heap holds $\{3,4,7\}$, the three smallest elements; its root
(maximum of these three) is $7$, the $3$rd smallest element overall.
Sorting $[3,4,7,10,15,20]$ confirms index $k-1=2$ holds $7$.

\subsection*{C++ implementation}

\begin{lstlisting}
#include <bits/stdc++.h>
using namespace std;

int kthSmallest(vector<int>& nums, int k) {
    priority_queue<int> maxHeap; // max-heap (default in C++)

    for (int x : nums) {
        maxHeap.push(x);
        if ((int)maxHeap.size() > k) {
            maxHeap.pop();
        }
    }
    return maxHeap.top();
}

int main() {
    vector<int> nums = {7, 10, 4, 3, 20, 15};
    cout << "3rd smallest: " << kthSmallest(nums, 3) << endl;
    return 0;
}
\end{lstlisting}

\begin{complexitybox}
\textbf{Time Complexity.}
\[
  O(n \log k),
\]
identical in structure to Section~\ref{sec:heap-kth-largest}.
\textbf{Space Complexity.} $O(k)$.
\end{complexitybox}

\begin{takeawaybox}
Once Kth Largest is internalized, Kth Smallest costs nothing extra to
derive -- it is the same algorithm with every comparison polarity
reversed. The general rule: a size-$k$ heap of the \emph{opposite}
polarity to what you are looking for always keeps the right boundary
element at the root.
\end{takeawaybox}

\section{Hand of Straights}
\label{sec:heap-hand-of-straights}

\problemstatement{We are given a hand of $n$ cards, $\textit{hand}[0..n-1]$
(integers, possibly with repeats), and an integer $\textit{groupSize}$.
Determine whether the hand can be rearranged into groups, each of exactly
$\textit{groupSize}$ cards with \textbf{consecutive} values (e.g.\
$[3,4,5]$ is a valid group of size $3$; $[3,4,6]$ is not). Return
\texttt{true} or \texttt{false}.}

\begin{patternbox}
\emph{``Always form the next group starting from the smallest remaining
card''} is the greedy insight, and \emph{``repeatedly find the current
smallest remaining distinct value''} is exactly what a min-heap (over
distinct values) is built for. Whenever a problem requires repeatedly
peeling off the globally smallest (or largest) remaining item from a
\emph{changing} collection -- not just once, but again and again as the
collection shrinks -- that repeated ``smallest of what's left'' query is
the heap keyword to watch for, distinguishing this from a single-pass
greedy that only needs one sorted scan.
\end{patternbox}

\subsection*{Understanding the problem carefully}

First, if $n$ is not divisible by $\textit{groupSize}$, the answer is
immediately \texttt{false} -- the cards cannot be partitioned evenly.
Otherwise, count the occurrences of each distinct card value (a hash map),
and observe a key greedy fact: whichever card value is currently the
\textbf{smallest} among all cards not yet placed into a group \emph{must}
be the starting card of some group -- it cannot be the second or later
card of any group, since no smaller card remains to precede it. So we are
forced to start a new group at the current minimum, every time.

\begin{ideabox}[Greedy Choice, Powered by a Min-Heap]
Maintain a min-heap of the \emph{distinct} card values currently present
in the hand (with their counts tracked separately in a hash map). Repeatedly: peek
the smallest remaining value $v$; this must start a new group. Check that
$v, v+1, \dots, v+\textit{groupSize}-1$ are all available (count $>0$ in
the hash map); if any is missing, return \texttt{false}. Otherwise,
decrement the count of each of these $\textit{groupSize}$ values; whenever
a value's count reaches $0$, remove it from consideration (pop it from the
heap, or simply leave it and skip over exhausted values when they resurface
at the top). Repeat until all cards are consumed.
\end{ideabox}

\subsection*{Why always starting at the current minimum is correct}

This is another exchange-argument greedy: suppose, for contradiction, that
some valid grouping does \emph{not} use the smallest remaining card $v$ as
the start of a group, but instead places it later in some group starting
at $v' < v$ -- but $v$ \emph{is} the minimum remaining, so no $v' < v$ can
exist among the remaining cards. Therefore $v$ cannot appear anywhere
except as the start of a group (it cannot be the second card of a group
starting at $v-1$, because $v-1$ is not present among the remaining
cards by assumption that $v$ is the minimum). So every valid grouping is
forced to start a group at $v$, and the greedy choice has no alternative
to consider in the first place -- it is not just \emph{a} good choice, it
is the \emph{only} choice consistent with feasibility.

\subsection*{Algorithm}

\begin{enumerate}
  \item If $n \bmod \textit{groupSize} \ne 0$: return \texttt{false}.
  \item Build a frequency map \textit{count} of card values; build a
        min-heap of the distinct values present.
  \item While the heap is not empty:
  \begin{enumerate}
    \item Pop values from the top of the heap whose count has reached
          $0$ (already fully used) and discard them.
    \item Peek the new minimum $v$ (if the heap is now empty, stop --
          all cards are used).
    \item For $j \gets 0$ to $\textit{groupSize}-1$: if
          $\textit{count}[v+j] = 0$, return \texttt{false}; otherwise
          decrement $\textit{count}[v+j]$.
  \end{enumerate}
  \item Return \texttt{true}.
\end{enumerate}

\begin{algorithm}[H]
\caption{Hand of Straights}
\begin{algorithmic}[1]
\Procedure{IsNStraightHand}{\textit{hand}, \textit{groupSize}}
  \If{$n \bmod \textit{groupSize} \ne 0$} \Return \textbf{false} \EndIf
  \State build frequency map \textit{count}; build min-heap of distinct values
  \While{\textit{heap} is not empty}
    \While{\textit{heap} not empty \textbf{and} $\textit{count}[\text{top}] = 0$}
      \State pop \textit{heap}
    \EndWhile
    \If{\textit{heap} is empty} \textbf{break} \EndIf
    \State $v \gets$ top of \textit{heap}
    \For{$j \gets 0$ to $\textit{groupSize}-1$}
      \If{$\textit{count}[v+j] = 0$} \Return \textbf{false} \EndIf
      \State $\textit{count}[v+j] \gets \textit{count}[v+j] - 1$
    \EndFor
  \EndWhile
  \State \Return \textbf{true}
\EndProcedure
\end{algorithmic}
\end{algorithm}

\subsection*{Dry run}

\dryrun{Take $\textit{hand} = [1,2,3,6,2,3,4,7,8]$,
$\textit{groupSize} = 3$.}

$n=9$, $9 \bmod 3 = 0$. Counts: $1{:}1, 2{:}2, 3{:}2, 4{:}1, 6{:}1, 7{:}1,
8{:}1$. Min-heap of distinct values: $\{1,2,3,4,6,7,8\}$.

\begin{itemize}
  \item Min is $1$. Group $[1,2,3]$: all counts $>0$. Decrement:
        $1{:}0, 2{:}1, 3{:}1$.
  \item Top of heap is $1$, but $\textit{count}[1]=0$ now; pop it. New
        min is $2$. Group $[2,3,4]$: all counts $>0$. Decrement:
        $2{:}0, 3{:}0, 4{:}0$.
  \item Pop $2,3,4$ as their counts reach $0$ at the top. New min is $6$.
        Group $[6,7,8]$: all counts $>0$ ($1$ each). Decrement to $0$
        each.
  \item Heap empties out as exhausted values are popped.
\end{itemize}

All cards consumed without failure, so the answer is \texttt{true}: the
hand splits into $[1,2,3]$, $[2,3,4]$, $[6,7,8]$.

\subsection*{C++ implementation}

\begin{lstlisting}
#include <bits/stdc++.h>
using namespace std;

bool isNStraightHand(vector<int>& hand, int groupSize) {
    int n = hand.size();
    if (n % groupSize != 0) return false;

    unordered_map<int,int> count;
    for (int x : hand) count[x]++;

    priority_queue<int, vector<int>, greater<int>> minHeap;
    for (auto& [val, c] : count) minHeap.push(val);

    while (!minHeap.empty()) {
        while (!minHeap.empty() && count[minHeap.top()] == 0) {
            minHeap.pop();
        }
        if (minHeap.empty()) break;

        int v = minHeap.top();
        for (int j = 0; j < groupSize; j++) {
            if (count[v + j] == 0) return false;
            count[v + j]--;
        }
    }
    return true;
}

int main() {
    vector<int> hand = {1, 2, 3, 6, 2, 3, 4, 7, 8};
    cout << (isNStraightHand(hand, 3) ? "true" : "false") << endl;
    return 0;
}
\end{lstlisting}

\begin{complexitybox}
\textbf{Time Complexity.} Building the frequency map and heap costs
$O(n \log n)$ in the worst case (up to $n$ distinct values, each pushed
once). The main loop performs at most $n/\textit{groupSize}$ group
formations, each doing $O(\textit{groupSize})$ work plus
$O(\log n)$ amortized heap-popping of exhausted values, giving an overall
time complexity of
\[
  O(n \log n).
\]
\textbf{Space Complexity.} $O(n)$ for the frequency map and heap in the
worst case (all distinct values).
\end{complexitybox}

\begin{takeawaybox}
When a greedy choice is not just \emph{good} but \emph{forced} (the
smallest remaining element has no alternative role to play), a min-heap
is the natural data structure to repeatedly surface that forced choice
efficiently, especially when -- unlike a single sorted pass -- elements
get consumed and removed from consideration in a pattern that does not
simply move left to right.
\end{takeawaybox}

\section{Merge K Sorted Lists}
\label{sec:heap-merge-k-lists}

\problemstatement{We are given $k$ linked lists, each individually sorted
in increasing order. Merge all of them into a single sorted linked list
and return it.}

\begin{patternbox}
\emph{``Merge $k$ sorted ...''} is one of the cleanest heap triggers in
this chapter: merging \emph{two} sorted lists is a familiar $O(n)$
two-pointer operation, but merging $k$ of them at once naturally wants a
data structure that can always hand you the smallest \emph{current head}
among all $k$ lists in $O(\log k)$ time -- exactly a min-heap holding one
candidate node from each list.
\end{patternbox}

\subsection*{Understanding the problem carefully}

At any point during the merge, the next node to append to the output is
the smallest value among the \emph{current heads} of the $k$ lists that
still have remaining nodes (every list is individually sorted, so only its
current head -- not any later node -- can possibly be the next-smallest
value globally). A naive approach would scan all $k$ current heads every
time to find the minimum, costing $O(k)$ per output node and $O(nk)$
overall (where $n$ is the total number of nodes). A min-heap reduces the
``find the minimum among $k$ candidates'' step to $O(\log k)$.

\begin{ideabox}[Greedy Choice, Powered by a Min-Heap]
Initialize a min-heap with the head node of each of the $k$ lists (if
non-null). Repeatedly: pop the smallest node from the heap, append it to
the output list, and if that node has a successor within its own original
list, push that successor onto the heap. Continue until the heap is
empty.
\end{ideabox}

\subsection*{Why popping the heap's minimum is always the correct next output node}

The heap, at every step, holds exactly one ``current head'' candidate per
list that still has unconsumed nodes -- and because each list is
individually sorted, the true global minimum of all \emph{remaining}
nodes across all lists must be one of these $k$ current heads (a node
deeper inside a list cannot be smaller than that same list's own current
head). So the heap's minimum at any moment is guaranteed to be the correct
next node for the merged output, and pushing that node's successor
immediately after popping it maintains the invariant that the heap always
holds the current head of every still-nonempty list.

\subsection*{Algorithm}

\begin{enumerate}
  \item Initialize a min-heap (ordered by node value) with the head of
        each of the $k$ lists, skipping any empty lists.
  \item Initialize an empty output list (with a dummy head for easy
        appending).
  \item While the heap is not empty:
  \begin{enumerate}
    \item Pop the node with the smallest value; append it to the output
          list.
    \item If this node has a \texttt{next} pointer (within its original
          list): push that next node onto the heap.
  \end{enumerate}
  \item Return the output list (skipping the dummy head).
\end{enumerate}

\begin{algorithm}[H]
\caption{Merge K Sorted Lists}
\begin{algorithmic}[1]
\Procedure{MergeKLists}{\textit{lists}}
  \State $\textit{heap} \gets$ empty min-heap (by node value)
  \For{each list $L$ in \textit{lists}}
    \If{$L$ is not empty} push $L$'s head onto \textit{heap} \EndIf
  \EndFor
  \State $\textit{dummy} \gets$ new empty node; $\textit{tail} \gets \textit{dummy}$
  \While{\textit{heap} is not empty}
    \State $\textit{node} \gets$ pop minimum from \textit{heap}
    \State $\textit{tail}.\textit{next} \gets \textit{node}$; $\textit{tail} \gets \textit{node}$
    \If{$\textit{node}.\textit{next} \ne \text{null}$}
      \State push $\textit{node}.\textit{next}$ onto \textit{heap}
    \EndIf
  \EndWhile
  \State \Return $\textit{dummy}.\textit{next}$
\EndProcedure
\end{algorithmic}
\end{algorithm}

\subsection*{Dry run}

\dryrun{Take three lists: $L_1 = 1\to4\to5$, $L_2=1\to3\to4$,
$L_3=2\to6$.}

\begin{itemize}
  \item Initialize heap with the three heads: $\{1_{(L_1)}, 1_{(L_2)},
        2_{(L_3)}\}$.
  \item Pop a $1$ (say from $L_1$): output $=1$. Push $L_1$'s next, $4$.
        Heap $=\{1_{(L_2)}, 2_{(L_3)}, 4_{(L_1)}\}$.
  \item Pop $1$ (from $L_2$): output $=1,1$. Push $L_2$'s next, $3$.
        Heap $=\{2_{(L_3)}, 3_{(L_2)}, 4_{(L_1)}\}$.
  \item Pop $2$ (from $L_3$): output $=1,1,2$. Push $L_3$'s next, $6$.
        Heap $=\{3_{(L_2)}, 4_{(L_1)}, 6_{(L_3)}\}$.
  \item Pop $3$ (from $L_2$): output $=1,1,2,3$. Push $L_2$'s next, $4$.
        Heap $=\{4_{(L_1)}, 4_{(L_2)}, 6_{(L_3)}\}$.
  \item Pop a $4$ (say from $L_1$): output $=1,1,2,3,4$. Push $L_1$'s
        next, $5$. Heap $=\{4_{(L_2)}, 5_{(L_1)}, 6_{(L_3)}\}$.
  \item Pop $4$ (from $L_2$): output $=1,1,2,3,4,4$. $L_2$ exhausted, push
        nothing. Heap $=\{5_{(L_1)}, 6_{(L_3)}\}$.
  \item Pop $5$: output $=\dots,5$. $L_1$ exhausted. Heap $=\{6_{(L_3)}\}$.
  \item Pop $6$: output $=\dots,6$. Heap empty.
\end{itemize}

Final merged list: $1,1,2,3,4,4,5,6$.

\subsection*{C++ implementation}

\begin{lstlisting}
#include <bits/stdc++.h>
using namespace std;

struct ListNode {
    int val;
    ListNode* next;
    ListNode(int v) : val(v), next(nullptr) {}
};

struct Compare {
    bool operator()(ListNode* a, ListNode* b) {
        return a->val > b->val; // min-heap by value
    }
};

ListNode* mergeKLists(vector<ListNode*>& lists) {
    priority_queue<ListNode*, vector<ListNode*>, Compare> minHeap;

    for (ListNode* node : lists) {
        if (node) minHeap.push(node);
    }

    ListNode dummy(0);
    ListNode* tail = &dummy;

    while (!minHeap.empty()) {
        ListNode* node = minHeap.top();
        minHeap.pop();

        tail->next = node;
        tail = node;

        if (node->next) minHeap.push(node->next);
    }
    return dummy.next;
}

int main() {
    ListNode* l1 = new ListNode(1);
    l1->next = new ListNode(4);
    l1->next->next = new ListNode(5);

    ListNode* l2 = new ListNode(1);
    l2->next = new ListNode(3);
    l2->next->next = new ListNode(4);

    ListNode* l3 = new ListNode(2);
    l3->next = new ListNode(6);

    vector<ListNode*> lists = {l1, l2, l3};
    ListNode* merged = mergeKLists(lists);

    for (ListNode* cur = merged; cur; cur = cur->next) {
        cout << cur->val << " ";
    }
    cout << endl;
    return 0;
}
\end{lstlisting}

\begin{complexitybox}
\textbf{Time Complexity.} The heap never holds more than $k$ nodes. Every
one of the $n$ total nodes (across all $k$ lists) is pushed once and
popped once, each costing $O(\log k)$, giving an overall time complexity
of
\[
  O(n \log k).
\]
\textbf{Space Complexity.} $O(k)$ for the heap (the output list reuses
the input nodes rather than allocating new ones).
\end{complexitybox}

\begin{takeawaybox}
``Merge $k$ sources, each individually sorted'' always reduces to ``track
one current candidate per source in a min-heap, pop the global minimum,
advance that source.'' This template applies identically whether the
sources are linked lists (this section), arrays, or streams, and is one of
the most directly reusable patterns in this chapter.
\end{takeawaybox}

\section{Replace Elements by Rank in the Array}
\label{sec:heap-replace-by-rank}

\problemstatement{Given an array $\textit{arr}$ of $n$ integers, replace
every element with its \textbf{rank} when the array's distinct values are
sorted in increasing order: the smallest value gets rank $1$, the next
distinct value gets rank $2$, and so on; equal elements receive the same
rank.}

\begin{patternbox}
\emph{``Replace each element by its rank in sorted order''} is, at its
core, a sorting problem -- but presenting it via a min-heap makes the
\textbf{``process elements in increasing order while remembering where
each one came from''} pattern explicit, which is the same underlying idea
used throughout this chapter whenever output order must differ from input
order. Whenever you need to process values in sorted order \emph{and}
write results back to their original positions, pair each value with its
original index before pushing it into the heap.
\end{patternbox}

\subsection*{Understanding the problem carefully}

Rank assignment requires processing values from smallest to largest while
remembering, for each value, which original array position it came from
(so the computed rank can be written back to the correct place). A
min-heap of $(\textit{value}, \textit{originalIndex})$ pairs, ordered by
value, gives exactly this controlled smallest-to-largest traversal.

\begin{ideabox}[Greedy Choice, Powered by a Min-Heap]
Push every $(\textit{value}, \textit{originalIndex})$ pair into a
min-heap ordered by value. Maintain a running rank counter, starting at
$0$, and a record of the previously popped value (initially undefined).
Repeatedly pop the smallest remaining pair: if its value differs from the
previously popped value (or this is the first pop), increment the rank
counter; write the current rank counter's value into the result array at
the popped pair's original index.
\end{ideabox}

\subsection*{Why incrementing the rank only on a value change is correct}

Equal elements must receive the same rank, and the rank must strictly
increase only when a genuinely new, larger distinct value is encountered.
Since the heap pops values in non-decreasing order, comparing each popped
value to the immediately previous one correctly detects exactly the
moments where a new distinct value begins -- there is no need to
pre-deduplicate the array, because consecutive equal pops are detected and
handled identically (no rank increment) regardless of how many times a
value repeats.

\subsection*{Algorithm}

\begin{enumerate}
  \item Push every $(\textit{arr}[i], i)$ pair onto a min-heap, ordered
        by value.
  \item Initialize $\textit{rank} \gets 0$, $\textit{prevValue} \gets$
        undefined.
  \item While the heap is not empty:
  \begin{enumerate}
    \item Pop $(\textit{value}, \textit{idx})$.
    \item If $\textit{value} \ne \textit{prevValue}$ (or this is the
          first pop): $\textit{rank} \gets \textit{rank}+1$;
          $\textit{prevValue} \gets \textit{value}$.
    \item $\textit{result}[\textit{idx}] \gets \textit{rank}$.
  \end{enumerate}
  \item Return \textit{result}.
\end{enumerate}

\begin{algorithm}[H]
\caption{Replace Elements by Rank}
\begin{algorithmic}[1]
\Procedure{ReplaceByRank}{\textit{arr}}
  \State $\textit{heap} \gets$ min-heap of $(\textit{arr}[i], i)$ for all $i$
  \State $\textit{rank} \gets 0$, $\textit{prevValue} \gets$ undefined
  \While{\textit{heap} is not empty}
    \State $(\textit{value}, \textit{idx}) \gets$ pop minimum
    \If{$\textit{value} \ne \textit{prevValue}$}
      \State $\textit{rank} \gets \textit{rank} + 1$
      \State $\textit{prevValue} \gets \textit{value}$
    \EndIf
    \State $\textit{result}[\textit{idx}] \gets \textit{rank}$
  \EndWhile
  \State \Return \textit{result}
\EndProcedure
\end{algorithmic}
\end{algorithm}

\subsection*{Dry run}

\dryrun{Take $\textit{arr} = [20, 15, 26, 2, 98, 6]$.}

Heap pairs: $(20,0),(15,1),(26,2),(2,3),(98,4),(6,5)$.

\begin{itemize}
  \item Pop $(2,3)$: new value. $\textit{rank}=1$.
        $\textit{result}[3]=1$.
  \item Pop $(6,5)$: new value. $\textit{rank}=2$.
        $\textit{result}[5]=2$.
  \item Pop $(15,1)$: new value. $\textit{rank}=3$.
        $\textit{result}[1]=3$.
  \item Pop $(20,0)$: new value. $\textit{rank}=4$.
        $\textit{result}[0]=4$.
  \item Pop $(26,2)$: new value. $\textit{rank}=5$.
        $\textit{result}[2]=5$.
  \item Pop $(98,4)$: new value. $\textit{rank}=6$.
        $\textit{result}[4]=6$.
\end{itemize}

Final result: $[4,3,5,1,6,2]$.

\subsection*{C++ implementation}

\begin{lstlisting}
#include <bits/stdc++.h>
using namespace std;

vector<int> replaceByRank(vector<int>& arr) {
    int n = arr.size();
    priority_queue<pair<int,int>, vector<pair<int,int>>,
                   greater<pair<int,int>>> minHeap;

    for (int i = 0; i < n; i++) {
        minHeap.push({arr[i], i});
    }

    vector<int> result(n);
    int rank = 0;
    bool first = true;
    int prevValue = -1;

    while (!minHeap.empty()) {
        auto [value, idx] = minHeap.top();
        minHeap.pop();

        if (first || value != prevValue) {
            rank++;
            prevValue = value;
            first = false;
        }
        result[idx] = rank;
    }
    return result;
}

int main() {
    vector<int> arr = {20, 15, 26, 2, 98, 6};
    vector<int> result = replaceByRank(arr);

    for (int x : result) cout << x << " ";
    cout << endl;
    return 0;
}
\end{lstlisting}

\begin{complexitybox}
\textbf{Time Complexity.} Pushing $n$ pairs costs $O(n \log n)$, and
popping all $n$ pairs costs another $O(n \log n)$, giving an overall time
complexity of
\[
  O(n \log n).
\]
\textbf{Space Complexity.} $O(n)$ for the heap and the result array.
\end{complexitybox}

\begin{takeawaybox}
Whenever a problem needs values processed in sorted order while results
must land back at their \emph{original} positions, push
$(\textit{value}, \textit{originalIndex})$ pairs together, rather than
sorting the values alone and losing track of where each one belongs. This
small habit -- carry the index along with the value -- recurs constantly
whenever sorting (or heap order) and output position diverge.
\end{takeawaybox}

\section{Task Scheduler}
\label{sec:heap-task-scheduler}

\problemstatement{We are given an array $\textit{tasks}$, where each
element is a task type (e.g.\ a character), and an integer $n$, the
\textbf{cooldown}: the same task type cannot be executed again until at
least $n$ units of time have passed since its last execution (the CPU may
sit \textbf{idle} for a unit of time if no eligible task is available).
Find the \textbf{minimum total time} needed to execute all tasks.}

\begin{patternbox}
\emph{``Always schedule whichever task type currently has the most
remaining occurrences (subject to a cooldown), repeatedly''} is a
perpetually-resorted greedy priority -- exactly the
\emph{``schedule/group with a cost or cooldown that always prefers the
largest remaining count first''} keyword pattern from
Section~\ref{sec:heap-intro}. Scheduling the most frequent task type first,
whenever it is eligible, spreads out the most demanding task type as early
and as evenly as possible, minimizing the idle time it would otherwise
force later.
\end{patternbox}

\subsection*{Understanding the problem carefully}

Intuitively, the task type with the highest remaining count is the most
constraining: if we delay it, we are more likely to be forced into idle
slots later, because its cooldown keeps recurring. So at every time step,
we should greedily execute whichever \emph{eligible} task type currently
has the largest remaining count. A max-heap keyed by remaining count gives
us this ``most frequent remaining task'' in $O(\log d)$ time, where $d$ is
the number of distinct task types (at most $26$ for lowercase letters, so
effectively a small constant).

\begin{ideabox}[Greedy Choice, Powered by a Max-Heap]
Count the frequency of each task type and push these counts into a
max-heap. Process time in chunks of $n+1$ consecutive units (one full
cooldown cycle): within each chunk, repeatedly pop the most frequent
remaining task type, decrement its count, and execute it in the next time
slot; if its count is still positive after decrementing, set it aside
(it cannot be reused within this same cooldown cycle) to be pushed back
into the heap only after the chunk completes. If both the heap and the
set-aside list are empty before filling all $n+1$ slots of the current
chunk, stop immediately (no trailing idle time is needed); otherwise, any
unfilled slots in the chunk count as idle time.
\end{ideabox}

\subsection*{Why processing in cooldown-sized chunks is correct}

Within a single chunk of $n+1$ consecutive time slots, no task type can
appear twice (that is exactly what the cooldown $n$ forbids -- a task
executed at the start of a chunk cannot be executed again until at least
$n$ slots later, which is precisely the boundary of the next chunk). So a
chunk can hold at most one execution of each distinct task type, and
greedily filling each chunk with the currently most-frequent remaining
task types (up to $\min(d, n+1)$ of them, where $d$ is how many distinct
types still have remaining work) never wastes an opportunity: any
slot left empty in a chunk after all distinct remaining types have been
used \emph{must} be idle, because no task is eligible. Choosing the
\emph{most frequent} remaining types to fill each chunk (rather than any
other selection) is what minimizes the number of future chunks that will
be forced to leave slots idle, since it depletes the most constraining
task types fastest, evening out the workload as quickly as possible.

\subsection*{Algorithm}

\begin{enumerate}
  \item Count the frequency of each task type; push all counts into a
        max-heap.
  \item Initialize $\textit{time} \gets 0$.
  \item While the heap is not empty:
  \begin{enumerate}
    \item $\textit{tempList} \gets$ empty list (tasks used this cycle,
          temporarily held out of the heap).
    \item For $i \gets 0$ to $n$ ($n+1$ slots in this cycle):
    \begin{enumerate}
      \item If the heap is not empty: pop the largest count, decrement
            it by $1$; if still $>0$, append it to $\textit{tempList}$.
      \item $\textit{time} \gets \textit{time}+1$.
      \item If the heap is empty and $\textit{tempList}$ is empty: break
            out of this inner loop early (no trailing idle needed).
    \end{enumerate}
    \item Push every count in $\textit{tempList}$ back onto the heap.
  \end{enumerate}
  \item Return \textit{time}.
\end{enumerate}

\begin{algorithm}[H]
\caption{Task Scheduler}
\begin{algorithmic}[1]
\Procedure{LeastInterval}{\textit{tasks}, $n$}
  \State count frequencies; push all into max-heap \textit{heap}
  \State $\textit{time} \gets 0$
  \While{\textit{heap} is not empty}
    \State $\textit{temp} \gets$ empty list
    \For{$i \gets 0$ to $n$}
      \If{\textit{heap} is not empty}
        \State $c \gets$ pop maximum; $c \gets c - 1$
        \If{$c > 0$} append $c$ to \textit{temp} \EndIf
      \EndIf
      \State $\textit{time} \gets \textit{time} + 1$
      \If{\textit{heap} is empty \textbf{and} \textit{temp} is empty}
        \State \textbf{break}
      \EndIf
    \EndFor
    \State push every count in \textit{temp} back onto \textit{heap}
  \EndWhile
  \State \Return \textit{time}
\EndProcedure
\end{algorithmic}
\end{algorithm}

\subsection*{Dry run}

\dryrun{Take $\textit{tasks} = [A,A,A,B,B,B]$, $n = 2$.}

Counts: $A{:}3, B{:}3$. Heap $=\{3_A, 3_B\}$.

\begin{itemize}
  \item \textbf{Cycle 1} ($3$ slots, since $n+1=3$): pop $A$ (or $B$;
        ties broken arbitrarily), decrement to $2$, set aside. Slot
        used, $\textit{time}=1$. Pop $B$, decrement to $2$, set aside.
        Slot used, $\textit{time}=2$. Heap now empty, temp has
        $\{2_A, 2_B\}$ but temp is not empty, so we still fill the third
        slot: heap is empty, so this slot is idle. $\textit{time}=3$.
        Push $2_A, 2_B$ back: heap $=\{2_A, 2_B\}$.
  \item \textbf{Cycle 2}: pop $A$ (decrement to $1$, set aside),
        $\textit{time}=4$. Pop $B$ (decrement to $1$, set aside),
        $\textit{time}=5$. Third slot: heap empty, temp not empty, idle.
        $\textit{time}=6$. Push $1_A, 1_B$ back.
  \item \textbf{Cycle 3}: pop $A$ (decrement to $0$, do not set aside),
        $\textit{time}=7$. Pop $B$ (decrement to $0$, do not set aside),
        $\textit{time}=8$. Now heap is empty and temp is empty: break
        early, no third idle slot needed.
\end{itemize}

Total time is $8$, matching the well-known schedule
$A,B,\text{idle},A,B,\text{idle},A,B$.

\subsection*{C++ implementation}

\begin{lstlisting}
#include <bits/stdc++.h>
using namespace std;

int leastInterval(vector<char>& tasks, int n) {
    unordered_map<char,int> freq;
    for (char t : tasks) freq[t]++;

    priority_queue<int> maxHeap;
    for (auto& [ch, c] : freq) maxHeap.push(c);

    int time = 0;
    while (!maxHeap.empty()) {
        vector<int> temp;
        int slots = n + 1;

        for (int i = 0; i < slots; i++) {
            if (!maxHeap.empty()) {
                int c = maxHeap.top() - 1;
                maxHeap.pop();
                if (c > 0) temp.push_back(c);
            }
            time++;
            if (maxHeap.empty() && temp.empty()) break;
        }
        for (int c : temp) maxHeap.push(c);
    }
    return time;
}

int main() {
    vector<char> tasks = {'A','A','A','B','B','B'};
    cout << "Minimum time: " << leastInterval(tasks, 2) << endl;
    return 0;
}
\end{lstlisting}

\begin{complexitybox}
\textbf{Time Complexity.} There are at most $d \le 26$ distinct task
types, so the heap has size $O(1)$ from a complexity standpoint, and the
total number of pop/push operations is bounded by the total task count
$N$ (each task is popped and re-pushed $O(1)$ times across all cycles in
the typical analysis), giving
\[
  O(N \log d) = O(N),
\]
effectively linear in the number of tasks.
\textbf{Space Complexity.} $O(d)$ for the frequency map and heap.
\end{complexitybox}

\begin{takeawaybox}
A closed-form formula also exists for this specific problem
($\textit{answer} = \max(N, (\textit{maxFreq}-1)(n+1) + \textit{numTasksWithMaxFreq})$),
but the heap-based simulation shown here generalizes better: it directly
extends to variants with per-task-type costs, partial cooldowns, or other
constraints where a clean closed form may not exist, while the
underlying ``always schedule the currently most-constraining option'' greedy
idea remains unchanged.
\end{takeawaybox}

\section{Design Twitter}
\label{sec:heap-design-twitter}

\problemstatement{Design a simplified Twitter supporting: $\textit{postTweet}
(\textit{userId}, \textit{tweetId})$ -- a user composes a new tweet;
$\textit{getNewsFeed}(\textit{userId})$ -- return the $10$ most recent
tweet IDs from the given user and everyone they follow, most recent
first; $\textit{follow}(\textit{followerId}, \textit{followeeId})$; and
$\textit{unfollow}(\textit{followerId}, \textit{followeeId})$.}

\begin{patternbox}
\emph{``Most recent $k$ items, drawn from several independently
time-ordered sources''} is exactly Merge K Sorted Lists
(Section~\ref{sec:heap-merge-k-lists}) combined with the Kth Largest
size-bounded heap idea (Section~\ref{sec:heap-kth-largest}), applied
together: each followed user's tweet history is individually
time-ordered (most recent last, or equivalently, sorted by a timestamp),
and we need the top $10$ across all of them -- a direct fusion of two
patterns already established in this chapter, now applied inside a larger
system-design problem.
\end{patternbox}

\subsection*{Understanding the problem carefully}

Each user's tweets, in posting order, already form an individually sorted
sequence (sorted by time posted). Getting a news feed means merging the
relevant followed users' tweet sequences and taking the top $10$ most
recent -- structurally identical to merging $k$ sorted lists and only
needing the first $10$ elements of the merge, rather than the full merge.
We assign a global, strictly increasing \textit{timestamp} to every tweet
at the moment it is posted (a simple counter, incremented on every
$\textit{postTweet}$ call across \emph{all} users), so that ``most
recent'' always corresponds to ``largest timestamp,'' comparable across
different users.

\begin{ideabox}[Data Structures]
\begin{itemize}
  \item A global counter \textit{timestamp}, incremented on every post.
  \item For each user, a list of $(\textit{timestamp}, \textit{tweetId})$
        pairs, in posting order (so the most recent tweet is always at
        the end).
  \item For each user, a set of followee user IDs (a user is always
        considered to follow themselves, for news-feed purposes, without
        needing an explicit \textit{follow} call).
\end{itemize}
\end{ideabox}

\begin{ideabox}[News Feed via a Bounded Max-Heap]
For $\textit{getNewsFeed}(\textit{userId})$: for each followee (including
the user themselves), look only at their \emph{most recent} tweet first
(the tail of their tweet list) -- exactly like the ``current head of each
sorted source'' idea from merging $k$ sorted lists, just read from the end
instead of the front, since within a user's own list, later postings have
larger timestamps. Push each followee's most recent tweet into a max-heap
keyed by timestamp. Repeatedly pop the global maximum, record its tweet
ID, and if that same followee has an earlier tweet remaining (walking
backward through their list), push that earlier tweet in its place. Stop
after popping $10$ tweets or exhausting the heap.
\end{ideabox}

\subsection*{Why this bounded heap pop is correct and efficient}

This is precisely the merge-$k$-sorted-lists technique, with two
adjustments: we traverse each followee's list from most recent to least
recent (so the heap always offers the next-best candidate, by timestamp,
in the right direction), and we stop after $10$ pops rather than draining
the entire merge, since the problem only ever wants the top $10$. Stopping
early is valid precisely because a max-heap's repeated-pop sequence
\emph{is} a fully sorted (by timestamp, descending) traversal of the
union of all followees' tweets -- truncating that sorted traversal at $10$
elements is exactly the $10$ most recent tweets, with no need to ever
materialize the full merge.

\subsection*{Algorithm}

\begin{enumerate}
  \item \textsc{PostTweet}($\textit{userId}, \textit{tweetId}$): append
        $(\textit{timestamp}{+}{+}, \textit{tweetId})$ to
        $\textit{userId}$'s tweet list.
  \item \textsc{Follow}/\textsc{Unfollow}: add/remove
        $\textit{followeeId}$ from $\textit{followerId}$'s followee set.
  \item \textsc{GetNewsFeed}($\textit{userId}$):
  \begin{enumerate}
    \item For each followee $f$ in $\textit{userId}$'s followee set
          (plus $\textit{userId}$ itself): if $f$ has at least one
          tweet, push $(\textit{their last tweet's timestamp}, f,
          \textit{index of that tweet})$ onto a max-heap.
    \item Repeat up to $10$ times: pop the maximum; record its tweet ID;
          if the popped followee has an earlier tweet at
          $\textit{index}-1$, push that one.
    \item Return the recorded tweet IDs, in the order popped (already
          most-recent-first).
  \end{enumerate}
\end{enumerate}

\subsection*{Dry run}

\dryrun{User $1$ posts tweets $5$ (at timestamp $1$) and $3$ (at
timestamp $2$). User $2$ posts tweet $101$ (at timestamp $3$). User $1$
follows User $2$. \textsc{GetNewsFeed}($1$) is called.}

User $1$'s followees (including self): $\{1, 2\}$. Most recent tweet of
user $1$: $(2, 3)$ -- timestamp $2$, tweet $3$. Most recent tweet of user
$2$: $(3, 101)$ -- timestamp $3$, tweet $101$. Heap initialized with both.

\begin{itemize}
  \item Pop max: $(3, 101)$ from user $2$. Record tweet $101$. User $2$
        has no earlier tweet; nothing pushed back.
  \item Pop max: $(2, 3)$ from user $1$. Record tweet $3$. User $1$ has
        an earlier tweet at timestamp $1$ (tweet $5$); push $(1, 5)$.
  \item Pop max: $(1, 5)$ from user $1$. Record tweet $5$. No earlier
        tweet from user $1$.
  \item Heap empty; fewer than $10$ tweets total, so we stop here anyway.
\end{itemize}

The news feed is $[101, 3, 5]$, most recent first -- correct.

\subsection*{C++ implementation}

\begin{lstlisting}
#include <bits/stdc++.h>
using namespace std;

class Twitter {
    int timestamp = 0;
    unordered_map<int, vector<pair<int,int>>> tweets; // userId -> (time, tweetId)
    unordered_map<int, unordered_set<int>> following;  // userId -> followees

public:
    void postTweet(int userId, int tweetId) {
        tweets[userId].push_back({timestamp++, tweetId});
    }

    vector<int> getNewsFeed(int userId) {
        // (timestamp, userId, index into that user's tweet list)
        priority_queue<tuple<int,int,int>> maxHeap;

        auto consider = [&](int uid) {
            if (!tweets[uid].empty()) {
                int idx = tweets[uid].size() - 1;
                maxHeap.push({tweets[uid][idx].first, uid, idx});
            }
        };

        consider(userId);
        for (int f : following[userId]) consider(f);

        vector<int> result;
        while (!maxHeap.empty() && (int)result.size() < 10) {
            auto [time, uid, idx] = maxHeap.top();
            maxHeap.pop();
            result.push_back(tweets[uid][idx].second);

            if (idx > 0) {
                maxHeap.push({tweets[uid][idx-1].first, uid, idx-1});
            }
        }
        return result;
    }

    void follow(int followerId, int followeeId) {
        if (followerId != followeeId) following[followerId].insert(followeeId);
    }

    void unfollow(int followerId, int followeeId) {
        following[followerId].erase(followeeId);
    }
};

int main() {
    Twitter tw;
    tw.postTweet(1, 5);
    tw.postTweet(1, 3);
    tw.postTweet(2, 101);
    tw.follow(1, 2);

    for (int id : tw.getNewsFeed(1)) cout << id << " ";
    cout << endl;
    return 0;
}
\end{lstlisting}

\begin{complexitybox}
\textbf{Time Complexity.} $\textit{postTweet}$, $\textit{follow}$, and
$\textit{unfollow}$ are $O(1)$ amortized. $\textit{getNewsFeed}$ pushes at
most $f$ initial candidates (one per followee, $f \le$ number followed)
and performs at most $10$ pop/push rounds, each $O(\log f)$, giving
\[
  O(f \log f)
\]
per call, independent of the total number of tweets ever posted.
\textbf{Space Complexity.} $O(T)$ to store all $T$ tweets ever posted,
plus $O(f)$ per news feed query for the heap.
\end{complexitybox}

\begin{takeawaybox}
System-design problems frequently turn out to be direct compositions of
patterns you already know: this one is exactly ``merge $k$ sorted
sources'' (Section~\ref{sec:heap-merge-k-lists}) plus ``stop after the
top $k$'' (Section~\ref{sec:heap-kth-largest}), wrapped in bookkeeping for
followers and a global timestamp. Recognizing the composition is more
valuable than treating it as an unfamiliar new problem.
\end{takeawaybox}

\section{Find Median from a Data Stream}
\label{sec:heap-median-stream}

\problemstatement{Design a structure supporting $\textit{addNum}(x)$,
which adds an integer to a growing stream, and $\textit{findMedian}()$,
which returns the median of all numbers added so far, at any point,
efficiently after every insertion.}

\begin{patternbox}
\emph{``Running median, computed after every new element of a stream''}
is the canonical trigger for the \textbf{two-heap} pattern flagged in
Section~\ref{sec:heap-intro}: a median needs access to both ``the largest
of the lower half'' and ``the smallest of the upper half'' of the data
seen so far, and no single heap can offer both -- so we keep two heaps of
opposite polarity, each covering one half, and keep them balanced.
\end{patternbox}

\subsection*{Understanding the problem carefully}

If we maintain the invariant that all elements in a \textbf{max-heap}
called \textit{lower} are $\le$ all elements in a \textbf{min-heap} called
\textit{upper}, and that the two heaps' sizes differ by at most $1$, then:

\begin{itemize}
  \item If the total count is odd, the median is the top of whichever
        heap has one more element than the other.
  \item If the total count is even, the median is the average of
        $\textit{lower}$'s top (the largest of the lower half) and
        $\textit{upper}$'s top (the smallest of the upper half).
\end{itemize}

The entire problem reduces to maintaining this ``balanced split into a
max-heap lower half and a min-heap upper half'' invariant after every
insertion.

\begin{ideabox}[Insertion Rule]
On $\textit{addNum}(x)$: push $x$ into $\textit{lower}$ (the max-heap).
Then, to preserve the ``everything in \textit{lower} is $\le$ everything in
\textit{upper}'' invariant, pop $\textit{lower}$'s maximum and push it into
$\textit{upper}$ (this guarantees the largest element currently in
\textit{lower} is never larger than anything already confirmed to belong
in the upper half). Finally, to preserve the size-balance invariant: if
$\textit{upper}$ now has more elements than $\textit{lower}$, pop
$\textit{upper}$'s minimum and push it back into $\textit{lower}$.
\end{ideabox}

\subsection*{Why pushing into \textit{lower} first, then shuffling, is correct}

Always inserting $x$ into $\textit{lower}$ first (regardless of whether
$x$ is actually a ``low'' or ``high'' value) and then immediately moving
$\textit{lower}$'s current maximum into $\textit{upper}$ is a clean way to
guarantee correctness without an explicit case split on where $x$ belongs:
after this shuffle, $\textit{upper}$'s new minimum is exactly
$\min(x, \text{old } \textit{upper}\text{'s minimum})$, and one can verify
that whichever of these is smaller is correctly placed at the boundary,
while $\textit{lower}$'s new maximum is whatever was second-largest among
$\{x\} \cup \textit{lower}\text{'s old elements}$, preserving
$\textit{lower} \le \textit{upper}$ throughout. The subsequent
size-rebalancing step (moving one element back if \textit{upper} becomes
too large) maintains the size-difference-at-most-$1$ invariant, which is
exactly what makes the parity-based median formula above always correct.

\subsection*{Algorithm}

\begin{enumerate}
  \item \textsc{AddNum}($x$):
  \begin{enumerate}
    \item Push $x$ onto $\textit{lower}$ (max-heap).
    \item Pop $\textit{lower}$'s max; push it onto $\textit{upper}$
          (min-heap).
    \item If $|\textit{upper}| > |\textit{lower}|$: pop $\textit{upper}$'s
          min; push it back onto $\textit{lower}$.
  \end{enumerate}
  \item \textsc{FindMedian}():
  \begin{enumerate}
    \item If $|\textit{lower}| > |\textit{upper}|$: return
          $\textit{lower}$'s top.
    \item Else: return $(\textit{lower}\text{'s top} +
          \textit{upper}\text{'s top}) / 2$.
  \end{enumerate}
\end{enumerate}

\begin{algorithm}[H]
\caption{Median Maintainer (Two Heaps)}
\begin{algorithmic}[1]
\Procedure{AddNum}{$x$}
  \State push $x$ onto \textit{lower} (max-heap)
  \State push (pop \textit{lower}) onto \textit{upper} (min-heap)
  \If{$|\textit{upper}| > |\textit{lower}|$}
    \State push (pop \textit{upper}) onto \textit{lower}
  \EndIf
\EndProcedure
\Procedure{FindMedian}{}
  \If{$|\textit{lower}| > |\textit{upper}|$} \Return top of \textit{lower}
  \Else \Return $(\text{top of } \textit{lower} + \text{top of } \textit{upper}) / 2$ \EndIf
\EndProcedure
\end{algorithmic}
\end{algorithm}

\subsection*{Dry run}

\dryrun{Add $5, 15, 1, 3$ in order, checking the median after each.}

\begin{itemize}
  \item Add $5$: push to \textit{lower} $=\{5\}$. Pop max ($5$), push to
        \textit{upper} $=\{5\}$. Now $|\textit{upper}|=1 >
        |\textit{lower}|=0$: pop \textit{upper}'s min ($5$), push to
        \textit{lower}. \textit{lower}$=\{5\}$, \textit{upper}$=\{\}$.
        Median: $|\textit{lower}|>|\textit{upper}|$, so median $=5$.
  \item Add $15$: push to \textit{lower}$=\{5,15\}$ (max-heap top $15$).
        Pop max ($15$), push to \textit{upper}$=\{15\}$.
        $|\textit{upper}|=1=|\textit{lower}|=1$, no rebalance needed.
        Median: equal sizes, $(5+15)/2=10$.
  \item Add $1$: push to \textit{lower}$=\{1,5\}$ (top $5$). Pop max
        ($5$), push to \textit{upper}$=\{5,15\}$ (top $5$).
        $|\textit{upper}|=2>|\textit{lower}|=1$: pop \textit{upper}'s min
        ($5$), push to \textit{lower}$=\{1,5\}$ (top $5$).
        \textit{upper}$=\{15\}$. Median: $|\textit{lower}|=2>
        |\textit{upper}|=1$, median $=5$.
  \item Add $3$: push to \textit{lower}$=\{1,3,5\}$ (top $5$). Pop max
        ($5$), push to \textit{upper}$=\{5,15\}$ (top $5$).
        $|\textit{upper}|=2=|\textit{lower}|=2$, no rebalance. Median:
        equal sizes, $(\text{top of lower}=3 + \text{top of upper}=5)/2
        =4$.
\end{itemize}

Checking by hand: after $[5,15,1,3]$, sorted is $[1,3,5,15]$, median
$=(3+5)/2=4$. Correct.

\subsection*{C++ implementation}

\begin{lstlisting}
#include <bits/stdc++.h>
using namespace std;

class MedianFinder {
    priority_queue<int> lower; // max-heap: lower half
    priority_queue<int, vector<int>, greater<int>> upper; // min-heap: upper half

public:
    void addNum(int x) {
        lower.push(x);
        upper.push(lower.top());
        lower.pop();

        if (upper.size() > lower.size()) {
            lower.push(upper.top());
            upper.pop();
        }
    }

    double findMedian() {
        if (lower.size() > upper.size()) {
            return lower.top();
        }
        return (lower.top() + upper.top()) / 2.0;
    }
};

int main() {
    MedianFinder mf;
    for (int x : {5, 15, 1, 3}) {
        mf.addNum(x);
        cout << "Median after adding " << x << ": "
             << mf.findMedian() << endl;
    }
    return 0;
}
\end{lstlisting}

\begin{complexitybox}
\textbf{Time Complexity.} Each $\textit{addNum}$ performs a constant
number of heap pushes and pops, each $O(\log n)$, giving
\[
  O(\log n) \text{ per insertion}, \qquad O(1) \text{ per median query}.
\]
\textbf{Space Complexity.} $O(n)$ to store all elements seen so far,
split between the two heaps.
\end{complexitybox}

\begin{takeawaybox}
The two-heap pattern -- a max-heap for ``everything below the median'' and
a min-heap for ``everything above,'' kept balanced in size -- is the
standard way to maintain order statistics (median, or more generally any
fixed percentile with a fixed split ratio) over a stream that only grows.
Whenever a running statistic needs comparison to \emph{both} a lower and
an upper neighborhood simultaneously, suspect this pattern before reaching
for a single heap.
\end{takeawaybox}

\section{Kth Largest Element in a Stream}
\label{sec:heap-kth-largest-stream}

\problemstatement{Design a class initialized with an integer $k$ and an
initial array $\textit{nums}$, supporting $\textit{add}(\textit{val})$:
add $\textit{val}$ to the running stream and return the $k$-th largest
element among \textbf{all} elements seen so far (including the initial
array and every value added).}

\begin{patternbox}
This is Section~\ref{sec:heap-kth-largest}'s ``Kth Largest Element in an
Array'' made \textbf{incremental}: instead of computing the answer once
for a fixed array, we must maintain it continuously as new elements
stream in, one at a time. The keyword shift from a single static query to
a class with a repeated $\textit{add}$ method is the signal that the
size-$k$ min-heap should be built \emph{once} and \emph{kept alive}
across calls, rather than recomputed from scratch on every query.
\end{patternbox}

\subsection*{Understanding the problem carefully}

The min-heap-of-size-$k$ technique from
Section~\ref{sec:heap-kth-largest} already processes elements one at a
time and remains correct at \emph{every intermediate point} during that
processing -- the heap, after processing any prefix of elements, always
holds the $k$ largest elements seen \emph{so far}, with the $k$-th largest
sitting at the root. This means no new derivation is needed: we simply
keep the heap object alive as part of the class's state, seed it with the
initial array, and apply the exact same push-then-conditionally-pop step
to each newly added value.

\begin{ideabox}[Reuse: Make the Heap Persistent]
On construction: build a min-heap from the initial $\textit{nums}$,
trimming it down to size $k$ using the same push-then-pop-if-oversized
rule as Algorithm~\ref{sec:heap-kth-largest}. On each $\textit{add}(val)$:
apply that same rule to $\textit{val}$ alone, then return the heap's
current root.
\end{ideabox}

\subsection*{Algorithm}

\begin{enumerate}
  \item \textsc{Constructor}($k$, \textit{nums}): for each $x$ in
        \textit{nums}, push $x$ onto the min-heap; if size exceeds $k$,
        pop the minimum.
  \item \textsc{Add}($\textit{val}$): push $\textit{val}$ onto the
        min-heap; if size exceeds $k$, pop the minimum; return the
        heap's current top.
\end{enumerate}

\begin{algorithm}[H]
\caption{Kth Largest Element in a Stream}
\begin{algorithmic}[1]
\Procedure{Constructor}{$k$, \textit{nums}}
  \State $\textit{heap} \gets$ empty min-heap
  \For{each $x$ in \textit{nums}}
    \State push $x$ onto \textit{heap}
    \If{$|\textit{heap}| > k$} pop the minimum \EndIf
  \EndFor
\EndProcedure
\Procedure{Add}{\textit{val}}
  \State push \textit{val} onto \textit{heap}
  \If{$|\textit{heap}| > k$} pop the minimum \EndIf
  \State \Return top of \textit{heap}
\EndProcedure
\end{algorithmic}
\end{algorithm}

\subsection*{Dry run}

\dryrun{Take $k=3$, initial $\textit{nums}=[4,5,8,2]$, then call
$\textit{add}(3)$, $\textit{add}(5)$, $\textit{add}(10)$.}

\begin{itemize}
  \item Construct: push $4,5,8,2$ one at a time, trimming to size $3$
        each time it exceeds $3$. After all four: heap holds the $3$
        largest of $\{4,5,8,2\}$, which are $\{4,5,8\}$ (root $=4$).
  \item $\textit{add}(3)$: push $3$; heap $=\{3,4,5,8\}$, size $4>3$:
        pop min ($3$). Heap $=\{4,5,8\}$. Return $4$.
  \item $\textit{add}(5)$: push $5$; heap $=\{4,5,5,8\}$, size $4>3$:
        pop min ($4$). Heap $=\{5,5,8\}$. Return $5$.
  \item $\textit{add}(10)$: push $10$; heap $=\{5,5,8,10\}$, size $4>3$:
        pop min ($5$). Heap $=\{5,8,10\}$. Return $5$.
\end{itemize}

\subsection*{C++ implementation}

\begin{lstlisting}
#include <bits/stdc++.h>
using namespace std;

class KthLargest {
    int k;
    priority_queue<int, vector<int>, greater<int>> minHeap;

public:
    KthLargest(int k, vector<int>& nums) : k(k) {
        for (int x : nums) {
            minHeap.push(x);
            if ((int)minHeap.size() > k) minHeap.pop();
        }
    }

    int add(int val) {
        minHeap.push(val);
        if ((int)minHeap.size() > k) minHeap.pop();
        return minHeap.top();
    }
};

int main() {
    vector<int> nums = {4, 5, 8, 2};
    KthLargest kl(3, nums);

    cout << kl.add(3) << endl;
    cout << kl.add(5) << endl;
    cout << kl.add(10) << endl;
    return 0;
}
\end{lstlisting}

\begin{complexitybox}
\textbf{Time Complexity.} Construction costs $O(n \log k)$ for the
initial $n$ elements. Each subsequent $\textit{add}$ call costs
$O(\log k)$.
\textbf{Space Complexity.} $O(k)$, since the heap never exceeds size $k$.
\end{complexitybox}

\begin{takeawaybox}
Many ``streaming'' versions of array problems require no new algorithmic
idea at all -- only recognizing that the data structure from the static
version (here, a size-$k$ min-heap) is already incremental by
construction, and simply needs to be kept alive as persistent state
rather than rebuilt on every query.
\end{takeawaybox}

\section{Maximum Sum Combination}
\label{sec:heap-max-sum-combination}

\problemstatement{We are given two arrays $A$ and $B$, each of size $N$,
and an integer $K$. Consider every possible pair sum $A[i]+B[j]$ for
$0 \le i,j < N$ (there are $N^2$ such pairs). Find the $K$ \textbf{largest}
pair sums, in descending order, \textbf{without} computing all $N^2$ sums.}

\begin{patternbox}
\emph{``$K$ largest sums formed by combining one element from each of two
(sorted) collections''} is a generalization of Kth Largest
(Section~\ref{sec:heap-kth-largest}) to a combinatorially large candidate
space: there are $N^2$ possible pairs, far too many to generate and sort
when $K \ll N^2$. The keyword to notice is that the candidate sums have
\textbf{internal structure} (they come from a grid of sorted index pairs),
which a heap can exploit to generate only the $K$ largest sums directly,
without ever materializing the rest.
\end{patternbox}

\subsection*{Understanding the problem carefully}

Sort both $A$ and $B$ in \textbf{descending} order first. Now,
$A[0]+B[0]$ is guaranteed to be the single largest possible sum (the
largest element of each array, paired together). Crucially, once we know
$A[0]+B[0]$ is the maximum, the \emph{next}-largest sum must come from
either decreasing the $A$-index to $1$ (pairing $A[1]$ with $B[0]$) or
decreasing the $B$-index to $1$ (pairing $A[0]$ with $B[1]$) -- no other
pair can exceed both of these, since both arrays are sorted descending,
and any pair using indices further from $(0,0)$ in either coordinate can
only be smaller than at least one of these two candidates.

\begin{ideabox}[Greedy Choice, Powered by a Max-Heap over Index Pairs]
Maintain a max-heap of candidate \emph{index pairs} $(i,j)$, keyed by
$A[i]+B[j]$, along with a \textit{visited} set to avoid pushing the same
pair twice. Start with $(0,0)$ in the heap. Repeat $K$ times: pop the pair
with the largest sum, record that sum, then push $(i+1,j)$ and $(i,j+1)$
(if within bounds and not already visited) as the next candidates that
could potentially be the next-largest sum.
\end{ideabox}

\subsection*{Why only two neighbors need to be considered at each step}

This is the same idea as merging $k$ sorted lists
(Section~\ref{sec:heap-merge-k-lists}), generalized to a 2D grid of sorted
candidates rather than $k$ separate 1D lists: because both $A$ and $B$ are
sorted descending, the sum $A[i]+B[j]$ is non-increasing as $i$ or $j$
increases (holding the other fixed) -- so the grid of all possible sums is
itself ``sorted'' along both axes, exactly like the row/column-sorted
matrices in Sections~\ref{sec:bs-search-2d-matrix}--\ref{sec:bs-matrix-median}.
Once a pair $(i,j)$ has been popped as the current maximum, the only pairs
that could possibly be the \emph{next}-largest, among those not yet
considered, are its immediate neighbors $(i+1,j)$ and $(i,j+1)$ -- every
other still-uncovered pair is reachable only by passing through one of
these two first, so they are never skipped over by this expansion rule.

\subsection*{Algorithm}

\begin{enumerate}
  \item Sort $A$ and $B$ in descending order.
  \item Initialize a max-heap (by sum) with the single pair $(0,0)$;
        mark $(0,0)$ visited.
  \item Repeat $K$ times:
  \begin{enumerate}
    \item Pop $(i,j)$ with the largest sum $A[i]+B[j]$; record this sum.
    \item If $i+1 < N$ and $(i+1,j)$ not visited: push it; mark
          visited.
    \item If $j+1 < N$ and $(i,j+1)$ not visited: push it; mark
          visited.
  \end{enumerate}
  \item Return the $K$ recorded sums.
\end{enumerate}

\begin{algorithm}[H]
\caption{Maximum Sum Combination}
\begin{algorithmic}[1]
\Procedure{MaxSumCombination}{$A, B, K$}
  \State sort $A, B$ descending
  \State $\textit{heap} \gets$ max-heap with $(A[0]+B[0], 0, 0)$
  \State mark $(0,0)$ visited
  \State $\textit{result} \gets$ empty list
  \For{$\textit{count} \gets 1$ to $K$}
    \State $(\textit{sum}, i, j) \gets$ pop maximum
    \State append \textit{sum} to \textit{result}
    \If{$i+1<N$ \textbf{and} $(i{+}1,j)$ not visited}
      \State push $(A[i{+}1]+B[j], i{+}1, j)$; mark visited
    \EndIf
    \If{$j+1<N$ \textbf{and} $(i,j{+}1)$ not visited}
      \State push $(A[i]+B[j{+}1], i, j{+}1)$; mark visited
    \EndIf
  \EndFor
  \State \Return \textit{result}
\EndProcedure
\end{algorithmic}
\end{algorithm}

\subsection*{Dry run}

\dryrun{Take $A = [3,2]$, $B = [1,4]$, $K = 2$. Sorted descending:
$A=[3,2]$, $B=[4,1]$.}

\begin{itemize}
  \item Start: heap $=\{(7,0,0)\}$ (sum $A[0]+B[0]=3+4=7$). Visited
        $=\{(0,0)\}$.
  \item Pop $(7,0,0)$: record $7$. Push $(1,0)$: sum
        $A[1]+B[0]=2+4=6$. Push $(0,1)$: sum $A[0]+B[1]=3+1=4$. Heap
        $=\{6,4\}$. Visited $=\{(0,0),(1,0),(0,1)\}$.
  \item Pop $(6,1,0)$: record $6$. $K=2$ reached; stop.
\end{itemize}

The top $2$ sums are $[7,6]$. Checking by hand: all four pair sums are
$3{+}4=7$, $3{+}1=4$, $2{+}4=6$, $2{+}1=3$; sorted descending,
$[7,6,4,3]$, and the top $2$ are indeed $[7,6]$.

\subsection*{C++ implementation}

\begin{lstlisting}
#include <bits/stdc++.h>
using namespace std;

vector<int> maxSumCombination(vector<int> A, vector<int> B, int K) {
    int n = A.size();
    sort(A.rbegin(), A.rend());
    sort(B.rbegin(), B.rend());

    // (sum, i, j)
    priority_queue<tuple<int,int,int>> maxHeap;
    set<pair<int,int>> visited;

    maxHeap.push({A[0] + B[0], 0, 0});
    visited.insert({0, 0});

    vector<int> result;
    while (K-- > 0 && !maxHeap.empty()) {
        auto [sum, i, j] = maxHeap.top();
        maxHeap.pop();
        result.push_back(sum);

        if (i + 1 < n && !visited.count({i+1, j})) {
            maxHeap.push({A[i+1] + B[j], i+1, j});
            visited.insert({i+1, j});
        }
        if (j + 1 < n && !visited.count({i, j+1})) {
            maxHeap.push({A[i] + B[j+1], i, j+1});
            visited.insert({i, j+1});
        }
    }
    return result;
}

int main() {
    vector<int> A = {3, 2};
    vector<int> B = {1, 4};
    for (int x : maxSumCombination(A, B, 2)) cout << x << " ";
    cout << endl;
    return 0;
}
\end{lstlisting}

\begin{complexitybox}
\textbf{Time Complexity.} Sorting costs $O(N \log N)$. The main loop runs
$K$ times, each doing $O(1)$ pushes/pops of $O(\log K)$ cost (the heap
never holds more than $O(K)$ candidate pairs), giving an overall time
complexity of
\[
  O(N \log N + K \log K).
\]
\textbf{Space Complexity.} $O(K)$ for the heap and visited set.
\end{complexitybox}

\begin{takeawaybox}
``$K$ largest/smallest results from combining elements of two or more
sorted collections'' generalizes the merge-$k$-sorted-lists idea
(Section~\ref{sec:heap-merge-k-lists}) to a combinatorial candidate space:
instead of $k$ separate sequences, the candidates form an implicit sorted
grid, and a heap explores that grid one ``neighbor expansion'' at a time,
from the best candidate outward, generating exactly the $K$ best results
without ever enumerating the rest.
\end{takeawaybox}

\section{Minimum Cost to Connect Sticks}
\label{sec:heap-connect-sticks}

\problemstatement{We are given $n$ sticks with given lengths. We may
repeatedly choose any two sticks, combining them into a single stick whose
length is the sum of the two, at a \textbf{cost equal to that sum}. We
repeat until only one stick remains. Find the \textbf{minimum total cost}
to combine all sticks into one.}

\begin{patternbox}
\emph{``Repeatedly combine the two smallest remaining items, at a cost
equal to their combination,''} is a direct heap-driven greedy: every
combination's cost depends on the lengths being combined, and combining
\textbf{small} sticks early keeps their (now-larger) sum from being
re-combined and re-charged many times over as the process continues. This
is structurally the same idea behind Huffman coding (building an optimal
prefix-free code by always merging the two least-frequent symbols first),
though no prior knowledge of Huffman coding is needed to derive it here.
\end{patternbox}

\subsection*{Understanding the problem carefully}

Every stick's original length is ultimately added into the cost once for
every combination step it participates in afterward, directly or as part
of a larger merged stick. A stick combined \emph{early} (when it is still
short) participates in fewer total dollars of cost across the remaining
steps than a stick combined \emph{late} (after it has already grown large
through earlier merges). So, to minimize total cost, we always want the
currently smallest available sticks to be the ones combined next --
deferring large sticks (or large merged sums) to participate in as few
further combination steps as possible.

\begin{ideabox}[Greedy Choice, Powered by a Min-Heap]
Push every stick length into a min-heap. Repeat until only one element
remains: pop the two smallest sticks, add their sum to the running total
cost, and push that sum back into the heap as a new (merged) stick.
\end{ideabox}

\subsection*{Why always combining the two smallest sticks is optimal}

This is an exchange argument identical in spirit to Huffman coding's
optimality proof: suppose an optimal combination sequence does not combine
the two globally smallest sticks $x \le y$ at some point, but instead
combines one of them with some larger stick $z$ first. One can show that
swapping the order -- combining $x$ and $y$ first instead -- never
increases the total cost, because $x$ and $y$ being the smallest means
they contribute the least possible amount per combination step they are
involved in, and delaying them only ever increases the number of
combination steps (each adding their value again, transitively, through
ever-larger merged sums) that their lengths end up contributing to.
Repeatedly applying this swap argument shows that always combining the two
currently smallest available sticks, at every step, is never worse than
any alternative ordering.

\subsection*{Algorithm}

\begin{enumerate}
  \item Push every stick length into a min-heap.
  \item Initialize $\textit{totalCost} \gets 0$.
  \item While the heap has more than one element:
  \begin{enumerate}
    \item Pop the two smallest values $a, b$.
    \item $\textit{totalCost} \gets \textit{totalCost} + (a+b)$.
    \item Push $a+b$ back onto the heap.
  \end{enumerate}
  \item Return \textit{totalCost}.
\end{enumerate}

\begin{algorithm}[H]
\caption{Minimum Cost to Connect Sticks}
\begin{algorithmic}[1]
\Procedure{ConnectSticks}{\textit{sticks}}
  \State $\textit{heap} \gets$ min-heap of all stick lengths
  \State $\textit{totalCost} \gets 0$
  \While{$|\textit{heap}| > 1$}
    \State $a \gets$ pop minimum, $b \gets$ pop minimum
    \State $\textit{totalCost} \gets \textit{totalCost} + a + b$
    \State push $a+b$ onto \textit{heap}
  \EndWhile
  \State \Return \textit{totalCost}
\EndProcedure
\end{algorithmic}
\end{algorithm}

\subsection*{Dry run}

\dryrun{Take $\textit{sticks} = [2,4,3]$.}

Heap $=\{2,3,4\}$.

\begin{itemize}
  \item Pop $2,3$. Cost $+=5$. Push $5$. Heap $=\{4,5\}$.
        $\textit{totalCost}=5$.
  \item Pop $4,5$. Cost $+=9$. Push $9$. Heap $=\{9\}$.
        $\textit{totalCost}=14$.
  \item Only one element remains; stop.
\end{itemize}

The minimum total cost is $14$.

\subsection*{C++ implementation}

\begin{lstlisting}
#include <bits/stdc++.h>
using namespace std;

int connectSticks(vector<int>& sticks) {
    priority_queue<int, vector<int>, greater<int>> minHeap(
        sticks.begin(), sticks.end());

    int totalCost = 0;
    while (minHeap.size() > 1) {
        int a = minHeap.top(); minHeap.pop();
        int b = minHeap.top(); minHeap.pop();

        totalCost += (a + b);
        minHeap.push(a + b);
    }
    return totalCost;
}

int main() {
    vector<int> sticks = {2, 4, 3};
    cout << "Minimum cost: " << connectSticks(sticks) << endl;
    return 0;
}
\end{lstlisting}

\begin{complexitybox}
\textbf{Time Complexity.} Building the initial heap costs $O(n)$ (using
the bottom-up build-heap technique from
Section~\ref{sec:heap-min-to-max}, which most standard library
implementations, including \texttt{std::priority\_queue}'s range
constructor, use internally). Each of the $n-1$ combination steps performs
two pops and one push, each $O(\log n)$, giving an overall time complexity
of
\[
  O(n \log n).
\]
\textbf{Space Complexity.} $O(n)$ for the heap.
\end{complexitybox}

\begin{takeawaybox}
``Repeatedly merge the two cheapest/smallest remaining items, paying a
cost proportional to the merge'' is the Huffman-coding greedy pattern,
recognizable on sight once you have seen it: a min-heap that always
surfaces the two cheapest current items, merged and pushed back, builds an
optimal merge order in $O(n \log n)$ without ever needing to consider
merge orders explicitly.
\end{takeawaybox}

\section{Top K Frequent Elements}
\label{sec:heap-top-k-frequent}

\problemstatement{Given an array $\textit{nums}$ and an integer $k$,
return the $k$ elements that appear \textbf{most frequently} in the
array.}

\begin{patternbox}
\emph{``Top $k$ frequent''} is the size-$k$ heap technique from
Section~\ref{sec:heap-kth-largest}, applied not to the raw values
themselves but to a \textbf{derived key} -- here, frequency, computed once
via a hash map before the heap ever comes into play. Whenever a problem
asks for the top/bottom $k$ items ranked by something other than their
own raw value (frequency, distance, a custom score), first compute that
key for every item, then apply the familiar size-$k$ heap exactly as if
the key were the value.
\end{patternbox}

\subsection*{Understanding the problem carefully}

First, count the frequency of every distinct element using a hash map --
this converts the problem from ``top $k$ frequent elements'' into ``top
$k$ elements by a number we now have for each of them,'' which is exactly
the shape of Section~\ref{sec:heap-kth-largest}, just with $(\textit{key} =
\textit{frequency}, \textit{value} = \textit{element})$ pairs instead of
raw numbers.

\begin{ideabox}[Reuse, with a Derived Key]
Build a frequency map. Maintain a min-heap of size $k$, ordered by
frequency, holding $(\textit{frequency}, \textit{element})$ pairs.
Process each distinct element: push it; if the heap exceeds size $k$, pop
the pair with the smallest frequency (the current weakest member of the
top-$k$-frequent set). At the end, the heap holds exactly the $k$ most
frequent elements.
\end{ideabox}

\subsection*{Why the size-$k$ heap argument transfers unchanged}

The correctness argument from Section~\ref{sec:heap-kth-largest} never
actually depended on the heap being keyed by the element's own value --
it only depended on the key being a single comparable number that we want
to rank by. Replacing ``value'' with ``frequency of that value'' changes
nothing about why the size-$k$ min-heap correctly converges to the top-$k$
set: the same induction (the heap always holds the best $k$ keys seen so
far, by the same push-then-evict-the-worst logic) applies verbatim.

\subsection*{Algorithm}

\begin{enumerate}
  \item Build a frequency map of $\textit{nums}$.
  \item For each distinct $(\textit{element}, \textit{freq})$ pair:
  \begin{enumerate}
    \item Push $(\textit{freq}, \textit{element})$ onto a min-heap
          (ordered by \textit{freq}).
    \item If the heap's size exceeds $k$: pop the minimum.
  \end{enumerate}
  \item Return the elements remaining in the heap.
\end{enumerate}

\begin{algorithm}[H]
\caption{Top K Frequent Elements}
\begin{algorithmic}[1]
\Procedure{TopKFrequent}{\textit{nums}, $k$}
  \State build frequency map \textit{freq} of \textit{nums}
  \State $\textit{heap} \gets$ empty min-heap (by frequency)
  \For{each $(\textit{elem}, f)$ in \textit{freq}}
    \State push $(f, \textit{elem})$ onto \textit{heap}
    \If{$|\textit{heap}| > k$} pop the minimum \EndIf
  \EndFor
  \State \Return all elements remaining in \textit{heap}
\EndProcedure
\end{algorithmic}
\end{algorithm}

\subsection*{Dry run}

\dryrun{Take $\textit{nums} = [1,1,1,2,2,3]$, $k = 2$.}

Frequency map: $1{:}3, 2{:}2, 3{:}1$.

\begin{itemize}
  \item Push $(3,1)$: heap $=\{(3,1)\}$.
  \item Push $(2,2)$: heap $=\{(2,2),(3,1)\}$. Size $2 \le 2$.
  \item Push $(1,3)$: heap $=\{(1,3),(2,2),(3,1)\}$. Size $3>2$: pop min
        $(1,3)$. Heap $=\{(2,2),(3,1)\}$.
\end{itemize}

The final heap holds elements $\{2, 1\}$, the two most frequent values
(frequencies $2$ and $3$ respectively).

\subsection*{C++ implementation}

\begin{lstlisting}
#include <bits/stdc++.h>
using namespace std;

vector<int> topKFrequent(vector<int>& nums, int k) {
    unordered_map<int,int> freq;
    for (int x : nums) freq[x]++;

    priority_queue<pair<int,int>, vector<pair<int,int>>,
                   greater<pair<int,int>>> minHeap; // (freq, element)

    for (auto& [elem, f] : freq) {
        minHeap.push({f, elem});
        if ((int)minHeap.size() > k) minHeap.pop();
    }

    vector<int> result;
    while (!minHeap.empty()) {
        result.push_back(minHeap.top().second);
        minHeap.pop();
    }
    return result;
}

int main() {
    vector<int> nums = {1, 1, 1, 2, 2, 3};
    for (int x : topKFrequent(nums, 2)) cout << x << " ";
    cout << endl;
    return 0;
}
\end{lstlisting}

\begin{complexitybox}
\textbf{Time Complexity.} Building the frequency map costs $O(n)$.
Processing each of the $d \le n$ distinct elements costs $O(\log k)$,
giving an overall time complexity of
\[
  O(n + d \log k).
\]
\textbf{Space Complexity.} $O(d)$ for the frequency map plus $O(k)$ for
the heap.
\end{complexitybox}

\begin{takeawaybox}
This closing problem of the chapter is a deliberate reminder: nearly
every heap problem you will face is one of a small number of templates --
size-$k$ heap, merge-$k$-sources, two-balanced-heaps, or
always-take-the-best-remaining-greedy -- applied to whatever key the
problem actually cares about ranking by. Compute that key first (here,
frequency, via a hash map), and the heap logic that follows is almost
always one you have already written in this chapter.
\end{takeawaybox}

